% FST-DE_DarkEnergy_Skeleton_v1_de.tex
% Entwurf -- Dunkle Energie als residuale Vakuum-Freie-Energie
% Status: ENTWURF (v1)
% Datum: 2026-03-10

\documentclass[12pt,a4paper]{article}
\usepackage[utf8]{inputenc}
\usepackage[T1]{fontenc}
\usepackage{lmodern}
\usepackage{amsmath,amssymb,amsthm}
\usepackage{xcolor}
\usepackage{geometry}
\geometry{a4paper, margin=2.5cm}
\usepackage[ngerman]{babel}
\usepackage{hyperref}
\hypersetup{colorlinks=true, linkcolor=blue!70!black, citecolor=green!50!black, urlcolor=blue!70!black}

\theoremstyle{definition}
\newtheorem{definition}{Definition}[section]
\theoremstyle{plain}
\newtheorem{theorem}[definition]{Theorem}
\newtheorem{lemma}[definition]{Lemma}
\newtheorem{proposition}[definition]{Satz}
\newtheorem{corollary}[definition]{Korollar}
\newtheorem{axiom}[definition]{Axiom}
\theoremstyle{remark}
\newtheorem{remark}[definition]{Bemerkung}

% --- Eigene Befehle ---
\newcommand{\Leff}{\Lambda_{\mathrm{eff}}}
\newcommand{\LUV}{\Lambda_{\mathrm{UV}}}
\newcommand{\LIR}{\Lambda_{\mathrm{IR}}}
\newcommand{\rhov}{\rho_{\mathrm{vac}}}
\newcommand{\dd}{\mathrm{d}}

\title{\textbf{Dunkle Energie als residuale Vakuum-Freie-Energie:\\[4pt]
Eine thermodynamische Schranke f\"ur die kosmologische Konstante}}

\author{Lukas Geiger\\[2pt]
\small Unabh\"angiger Forscher, Bernau im Schwarzwald, Deutschland\\[2pt]
\small ORCID: \href{https://orcid.org/0009-0005-7296-1534}{0009-0005-7296-1534}%
\thanks{KI-Werkzeuge (Claude von Anthropic) wurden für mathematische Formalisierung, Literaturverifikation und redaktionelle Verfeinerung eingesetzt. Alle mathematischen Inhalte wurden vom Autor entwickelt, verifiziert und freigegeben.}}

\date{M\"arz 2026}

\begin{document}
\maketitle

% ============================================================
\begin{abstract}
Das Problem der kosmologischen Konstante geh\"ort zu den tiefsten
R\"atseln der theoretischen Physik: Naive Absch\"atzungen der
Vakuumenergiedichte aus der Quantenfeldtheorie \"uberschreiten den
beobachteten Wert um etwa 120 Gr\"o\ss enordnungen. Wir schlagen ein
thermodynamisches Variationsprinzip vor, das die effektive kosmologische
Konstante~$\Leff$ ohne Feinabstimmung einschr\"ankt. Ausgehend von einer
Freie-Energie-Funktion~$\Phi(\rhov)$, definiert \"uber
Versuchs-Vakuumenergiedichten und regularisiert zwischen einem
Ultraviolett-Cutoff~$\LUV$ und einem Infrarot-Cutoff, gegeben durch den
Hubble-Radius~$H_0^{-1}$, zeigen wir, dass das eindeutige Minimum
von~$\Phi$ eine residuale Vakuumenergiedichte liefert, die
wie~$\Leff \sim H_0^{2}/\ln(M_{\mathrm{Pl}}/H_0)$ skaliert. Die
resultierende numerische Absch\"atzung liegt innerhalb einer
Gr\"o\ss enordnung des Planck-2018-Wertes.  Durch Einbettung der
variationellen Schranke in einen Skalar-Tensor-Lagrangian
($R + \gamma R^{2}$ mit P\"oschl--Teller-Skalarfeld) leiten wir eine
tanh-S\"attigungsdynamik f\"ur den Dunkle-Energie-Dichteparameter, eine
effektive Zustandsgleichung im Thawing-Regime und eine falsifizierbare
Vorhersage f\"ur den gravitativen Slip-Parameter~$\eta = 5/4$ auf
sub-Compton-Skalen her.  Das Rahmenwerk adressiert auf nat\"urliche
Weise das Koinzidenzproblem, indem es die Vakuumenergie an die aktuelle
Expansionsrate koppelt.

\medskip\noindent
\textbf{Schl\"usselw\"orter:} Kosmologische Konstante, Dunkle Energie,
Vakuumenergie, thermodynamisches Variationsprinzip, Freie-Energie-Funktion,
Koinzidenzproblem, Skalar-Tensor-Gravitation, gravitativer Slip
\end{abstract}

% ============================================================
\section{Einleitung}\label{sec:intro}

Die kosmologische Konstante~$\Lambda$ wurde von Einstein 1917 eingef\"uhrt,
um ein statisches Universum zu erm\"oglichen, und sp\"ater nach der
Entdeckung der kosmischen Expansion wieder verworfen. Ihre moderne
Reinkarnation als Dunkle Energie -- der Treiber der beschleunigten
Expansion, die seit 1998 beobachtet
wird~\cite{Riess1998,Perlmutter1999} -- hat sie zur\"uck ins Zentrum
der fundamentalen Physik gebracht.

\subsection{Das Problem der kosmologischen Konstante}

In der Quantenfeldtheorie tr\"agt jedes Feld Nullpunktsenergie zum Vakuum
bei. Eine direkte Berechnung mit einem Planck-Skalen-Cutoff
ergibt~\cite{Weinberg1989,Zeldovich1968}
\begin{equation}\label{eq:qft-estimate}
  \rhov^{\mathrm{QFT}} \;\sim\; \frac{M_{\mathrm{Pl}}^{4}}{16\pi^{2}}
  \;\approx\; 10^{74}\;\mathrm{GeV}^{4},
\end{equation}
w\"ahrend die beobachtete Dunkle-Energie-Dichte
\begin{equation}\label{eq:obs}
  \rhov^{\mathrm{obs}} \;\approx\; 10^{-47}\;\mathrm{GeV}^{4}
\end{equation}
betr\"agt -- eine Diskrepanz von etwa~$10^{120}$.

\subsection{Bestehende Ans\"atze}

Verschiedene Strategien wurden verfolgt, um diese Diskrepanz aufzul\"osen
oder zu umgehen:
\begin{itemize}
  \item \textbf{Anthropische Selektion} in einer Landschaft von
        Vakua~\cite{BoussoPolchinski2000,Susskind2003}.
  \item \textbf{Quintessenz} -- ein langsam rollendes Skalarfeld, dessen
        potentielle Energie~$\Lambda$ nachahmt~\cite{PeeblesRatra2003}.
  \item \textbf{Modifizierte Gravitation} -- Ersetzung der Allgemeinen
        Relativit\"atstheorie auf kosmologischen Skalen ($f(R)$-Theorien,
        massive Gravitation).
  \item \textbf{Sequester-Mechanismen}, die die Vakuumenergie von der
        Kr\"ummung entkoppeln~\cite{KaloperPadilla2014}.
\end{itemize}
Jeder Ansatz l\"ost einige Aspekte, f\"uhrt jedoch neue freie Parameter
oder konzeptionelle Schwierigkeiten ein.

\subsection{Unser Ansatz}

Wir w\"ahlen einen anderen Weg. Anstatt die Vakuumenergie zu
kompensieren oder Selektionseffekte heranzuziehen, argumentieren wir,
dass die \emph{beobachtete} kosmologische Konstante der
thermodynamische Gleichgewichtswert eines Freie-Energie-Funktionals
ist, das \"uber Vakuumfluktuationen definiert ist. Das Funktional ist
nach unten durch die Infrarot-Geometrie der Raumzeit beschr\"ankt, und
sein eindeutiges Minimum erzeugt auf nat\"urliche Weise eine residuale
Energiedichte der Gr\"o\ss enordnung~$H_0^{2}/G$.

Diese variationelle Architektur realisiert die \textbf{,,Muster-B``-Logik (Flow-Selektion)} des FST-Rahmenwerks, wie sie in der \textbf{eichtheoretischen Neuformulierung der Riemannschen Hypothese} entwickelt wurde (siehe~\cite{GeigerRH2026c}). Ebenso wie die Stabilit\"at des arithmetischen Kerns durch eine Resolventen-Dominanz zweiter Ordnung gew\"ahrleistet wird, die die ,,physikalischen`` Nullstellen identifiziert, entsteht die kosmologische Konstante hier als der ,,Selektions-Eichparameter``, der das Vakuum gegen unkontrollierte Fluktuationen stabilisiert. Das Funktional $\Phi(\rho)$ ist das physikalische Analogon des renormierten Li-Funktionals im RH-Fall: Beide identifizieren einen eindeutigen, nicht-negativen Minimierer als Attraktor des komplexen Systems. Dunkle Energie ist somit der ,,Selektionsdruck``, der f\"ur gravitationale Stabilit\"at erforderlich ist.

% ============================================================
\section{Das Vakuum-Freie-Energie-Funktional}\label{sec:functional}

\subsection{Aufbau und Definitionen}

Betrachte einen r\"aumlich flachen Friedmann--Lema\^itre--Robertson--Walker-Hintergrund (FLRW)
mit Skalenfaktor~$a(t)$ und Hubble-Parameter~$H = \dot{a}/a$.
Wir definieren eine Freie-Energie-Funktion \"uber Versuchs-Vakuumenergiedichten~$\rho > 0$:

\begin{definition}[Vakuum-Freie-Energie-Funktion]\label{def:functional}
Seien $\LUV$ und $\LIR$ Ultraviolett- bzw.\ Infrarot-Impuls-Cutoffs.
Die \emph{Vakuum-Freie-Energie-Funktion} ist
\begin{equation}\label{eq:Phi}
  \Phi(\rho) \;=\; \int_{\LIR}^{\LUV}
    \left[\,
      \frac{1}{2}\,\omega(k)\;+\;
      T_{\mathrm{dS}}\,\rho\,\ln\!\left(\frac{\rho}{\rho_{\mathrm{Pl}}}\right)
      \;+\; V_{\mathrm{grav}}(\rho,\,k)
    \,\right]
    \frac{4\pi\,k^{2}}{(2\pi)^{3}}\;\dd k,
\end{equation}
wobei $\rho > 0$ eine Versuchs-Vakuumenergiedichte (ein skalarer Parameter, kein Feld) ist und:
\begin{itemize}
  \item $\omega(k) = \sqrt{k^{2} + m^{2}}$ die Modenfrequenz ist (summiert
        \"uber alle Spezies),
  \item $T_{\mathrm{dS}} = H/(2\pi)$ die de-Sitter-Temperatur (Gibbons--Hawking)
        des kosmologischen Horizonts ist,
  \item $\rho_{\mathrm{Pl}} = M_{\mathrm{Pl}}^{4}$ die
        Planck-Energiedichte ist,
  \item $V_{\mathrm{grav}}(\rho,k)$ die gravitative R\"uckwirkung der
        Vakuumenergie auf die Modenstruktur kodiert.
\end{itemize}
Der Entropieterm $\rho\,\ln(\rho/\rho_{\mathrm{Pl}})$ hat die thermodynamische Standardform $F = E - TS$ mit der Boltzmann-Entropiedichte $s(\rho) = -\rho\,\ln(\rho/\rho_{\mathrm{Pl}})$.
\end{definition}

\begin{remark}
Der Entropieterm $T_{\mathrm{dS}}\,\rho\,\ln(\rho/\rho_{\mathrm{Pl}})$
wirkt als thermodynamische Strafe.  F\"ur $\rho \to 0^{+}$ gilt
$\rho\,\ln(\rho/\rho_{\mathrm{Pl}}) \to 0^{-}$, und $V_{\mathrm{grav}}
\propto \rho^{2} \to 0$, sodass $\Phi$ auf die (endliche,
$\rho$-unabh\"angige) Nullpunktsenergie reduziert wird.  F\"ur
$\rho \to +\infty$ divergieren sowohl $\rho\,\ln\rho$ als auch
$\rho^{2}$, sodass $\Phi \to +\infty$.  Das Minimum liegt daher bei
einem Zwischenwert $\rho_{\ast} > 0$.
\end{remark}

\begin{remark}[Dimensionskonvention]
In nat\"urlichen Einheiten tragen die drei Terme unter dem
Integral in~\eqref{eq:Phi} verschiedene
Ingenieurdimensionen:
$[\omega(k)] = [\mathrm{Energie}]$,
$[T_{\mathrm{dS}}\,\rho\,\ln(\rho/\rho_{\mathrm{Pl}})]
= [\mathrm{Energie}]^{5}$ und
$[G\rho^{2}/k^{2}] = [\mathrm{Energie}]^{4}$.
Dies spiegelt den \emph{effektiven} Charakter des Funktionals
wider: Jeder Term wird implizit von dimensionslosen
Kopplungskonstanten der Gr\"o\ss enordnung Eins begleitet
(absorbiert in die Koeffizienten~$B$ und~$C$ im Beweis von
Theorem~\ref{thm:main}). Der physikalisch relevante Gehalt des
Variationsprinzips liegt in der $\rho$-Abh\"angigkeit jedes
Terms, die die Position des Minimums bestimmt. Die absolute
Gr\"o\ss e von~$\Phi$ ist nicht beobachtbar; nur die
Stationarit\"atsbedingung $\dd\Phi/\dd\rho = 0$ hat
physikalischen Gehalt, und diese Bedingung enth\"alt nur
Verh\"altnisse der Koeffizienten $B/C$, wobei die impliziten
Dimensionsfaktoren sich herausk\"urzen. Eine vollst\"andig
dimensionshomogene Formulierung---z.\,B.\ in Termen des
dimensionslosen Verh\"altnisses
$\rho/\rho_{\mathrm{Pl}}$---wird der vollst\"andigen Version
vorbehalten.
\end{remark}

\begin{remark}[Dimensionslose Formulierung]\label{rem:dimensionless}
Alle Terme im Funktional~\eqref{eq:Phi} werden dimensionslos, indem Dichten in Einheiten von $\rho_{\mathrm{Pl}} = c^5/(\hbar G^2)$, Wellenzahlen in Einheiten der Planck-L\"ange $\ell_{\mathrm{Pl}}$ und Energien in Einheiten von $M_{\mathrm{Pl}} c^2$ gemessen werden. In diesen nat\"urlichen Einheiten tragen der Entropiekoeffizient $T_{\mathrm{dS}}$, die Gravitationskopplung $G$ und die Modenenergie $\omega(k)$ s\"amtlich dimensionslose Vorfaktoren der Gr\"o\ss enordnung $O(1)$. Die im Text erw\"ahnten Kopplungskonstanten sind genau diese Planck-Einheiten-Normierungen.
\end{remark}

\subsection{Wahl der Cutoffs}

Der Ultraviolett-Cutoff wird an der reduzierten Planck-Skala angesetzt:
\begin{equation}
  \LUV \;=\; M_{\mathrm{Pl}} \;=\;
  \left(\frac{\hbar c}{8\pi G}\right)^{1/2}
  \;\approx\; 2.4 \times 10^{18}\;\mathrm{GeV}.
\end{equation}
Der Infrarot-Cutoff wird mit dem Hubble-Radius identifiziert:
\begin{equation}\label{eq:IR-cutoff}
  \LIR \;=\; H_0 \;\approx\; 1.5 \times 10^{-42}\;\mathrm{GeV}.
\end{equation}
Diese Identifikation hat Vorbilder in holographischen
Dunkle-Energie-Modellen~\cite{Li2004} und ist physikalisch durch die
Kausalstruktur des de-Sitter-Raums motiviert.

\subsection{Gravitativer R\"uckwirkungsterm}

Wir modellieren $V_{\mathrm{grav}}$ als die f\"uhrende
selbstgravitative Energiekosten der Aufrechterhaltung einer
Vakuumenergiedichte~$\rho$ auf der Skala~$k$:
\begin{equation}\label{eq:Vgrav}
  V_{\mathrm{grav}}(\rho,\,k) \;=\;
  \frac{8\pi G\,\rho^{2}}{k^{2}}.
\end{equation}
Diese quadratische Abh\"angigkeit von~$\rho$ stellt sicher, dass
$\Phi(\rho)$ f\"ur jede Mode strikt konvex ist, was f\"ur die
Eindeutigkeit des Minimums wesentlich ist.

\begin{lemma}[Minimalit\"at der R\"uckwirkungsform]\label{lem:minimality}
Sei $V[\rho, k]$ ein gravitatives R\"uckwirkungspotential, das folgende Bedingungen erf\"ullt:
\begin{itemize}
    \item[\textup{(M1)}] \textbf{Gravitationskopplung:} $V$ ist proportional zu genau einer Potenz der Newtonschen Konstante~$G$.
    \item[\textup{(M2)}] \textbf{Dimensionskonsistenz:} In nat\"urlichen Einheiten ($c = \hbar = 1$) gilt $[V] = [\mathrm{Energie}]^4$ (Energiedichte pro Mode).
    \item[\textup{(M3)}] \textbf{Minimale Nichtlinearit\"at:} $V$ ist das \emph{niedrigste} Monom in~$\rho$, das mit Selbstwechselwirkung vertr\"aglich ist (d.\,h.\ mindestens quadratisch in~$\rho$).
    \item[\textup{(M4)}] \textbf{IR-Dominanz:} $V$ w\"achst mit abnehmendem $k$ (infrarot-dominierte R\"uckwirkung).
\end{itemize}
Dann gilt, bis auf eine dimensionslose Konstante~$c_0 > 0$,
\[
V[\rho, k] \;=\; c_0\,\frac{G\,\rho^2}{k^2}.
\]
Die konventionelle Wahl $c_0 = 8\pi$ (einschlie\ss lich
Modenz\"ahlfaktoren) wird in~\eqref{eq:Vgrav} \"ubernommen.
\end{lemma}

\begin{proof}
Schreibe $V = G^\alpha \rho^\beta k^\gamma$ mit zu bestimmenden $\alpha, \beta, \gamma$. In nat\"urlichen Einheiten gilt $[G] = [\mathrm{Energie}]^{-2}$, $[\rho] = [\mathrm{Energie}]^4$, $[k] = [\mathrm{Energie}]^1$. Die Dimensionsbedingung $[V] = [\mathrm{Energie}]^4$ erfordert
\[
-2\alpha + 4\beta + \gamma = 4.
\]
Nach (M1) ist $\alpha = 1$, woraus $\gamma = 6 - 4\beta$ folgt. Nach~(M3) gilt $\beta \ge 2$. Die minimale Wahl $\beta = 2$ liefert $\gamma = -2$, also $V \sim G\,\rho^2 / k^2$. Der n\"achsth\"ohere Term ($\beta = 3$, $\gamma = -6$) ergibt $V \sim G\,\rho^3 / k^6$, was um den Faktor $\rho / k^4$ gegen\"uber dem f\"uhrenden Term unterdr\"uckt ist---vernachl\"assigbar im Regime $\rho \ll k^4$ (d.\,h.\ unterhalb der Planck-Dichte). Schlie\ss lich verlangt (M4) $\gamma < 0$, was durch $\gamma = -2$ erf\"ullt ist.

Der Wert von~$c_0$ wird durch Anpassung an die gravitative
Selbstenergie bestimmt. Die Newtonsche Potentialenergie einer
gleichf\"ormigen Kugel mit Dichte~$\rho$ und Radius~$R = k^{-1}$
betr\"agt $U = -\frac{16\pi^2}{15}\,G\,\rho^2\,R^5$, was eine
Selbstenergiedichte $u = |U|/(4\pi R^3/3) =
\frac{4\pi}{5}\,G\,\rho^2/k^2$ ergibt. Der $8\pi$-Koeffizient
in~\eqref{eq:Vgrav} stellt eine konventionelle Normierung dar, die
den Modenz\"ahlfaktor aus dem $k$-Raum-Integrationsma\ss{}
($4\pi k^2/(2\pi)^3$) einschlie\ss t. Die dimensionslose
Konstante~$c_0$ ist daher von der Gr\"o\ss enordnung
$O(4\pi)$--$O(8\pi)$; ihr pr\"aziser Wert beeinflusst die
parametrische Skalierung des Minimums nicht und kann in die
Renormierung von $\Lambda_{\mathrm{ren}}$ absorbiert werden.
\end{proof}

\begin{remark}
Lemma~\ref{lem:minimality} adressiert den zentralen Kritikpunkt, dass $V_{\mathrm{grav}}$ ein \emph{Ansatz} und keine Ableitung aus ersten Prinzipien ist. Obwohl das Lemma $V_{\mathrm{grav}}$ nicht aus den Einstein-Gleichungen herleitet, zeigt es, dass die funktionale Form~\eqref{eq:Vgrav} durch vier physikalisch transparente Axiome \emph{eindeutig bestimmt} ist (bis auf einen numerischen Koeffizienten). Jeder alternative R\"uckwirkungsterm w\"urde entweder die Dimensionskonsistenz verletzen, h\"ohere Gravitationskopplungen einf\"uhren ($G^2$-Terme, unterdr\"uckt um $M_{\mathrm{Pl}}^{-2}$) oder gegen\"uber dem f\"uhrenden Term subdominant sein.
\end{remark}

% ============================================================
\section{Hauptresultat}\label{sec:main}

\begin{theorem}[Thermodynamische Schranke f\"ur $\Leff$]\label{thm:main}
Unter den Voraussetzungen von Definition~\ref{def:functional} nimmt die
Funktion $\Phi(\rho)$ ein eindeutiges globales Minimum bei
$\rho = \rho_{\ast} > 0$ an. Die entsprechende effektive kosmologische
Konstante
\begin{equation}\label{eq:Leff}
  \Leff \;=\; \frac{8\pi G}{c^{4}}\;\rho_{\ast}
\end{equation}
erf\"ullt die Schranke
\begin{equation}\label{eq:bound}
  \Leff \;\sim\;
  \frac{H_0^{2}}{c^{2}}\;
  \ln\!\left(\frac{M_{\mathrm{Pl}}}{H_0}\right)^{-1}.
\end{equation}
\end{theorem}

\begin{proof}[Beweis (Skizze)]
Schreibe $\Phi(\rho) = A + B\,f(\rho) + C\,\rho^{2}$, wobei
$A = \int_{\LIR}^{\LUV} \tfrac{1}{2}\omega(k)\,\dd\mu(k)$ die
($\rho$-unabh\"angige) Nullpunktsenergie ist,
$B = T_{\mathrm{dS}}\int_{\LIR}^{\LUV}\dd\mu(k) > 0$ mit
$\dd\mu(k) = 4\pi k^{2}/(2\pi)^{3}\,\dd k$,
$f(\rho) = \rho\,\ln(\rho/\rho_{\mathrm{Pl}})$ und
$C = 8\pi G\int_{\LIR}^{\LUV} k^{-2}\,\dd\mu(k) > 0$.

\medskip\noindent
\textbf{Schritt~1: Existenz.}\;
F\"ur $\rho \to 0^{+}$: $f(\rho) = \rho\ln(\rho/\rho_{\mathrm{Pl}}) \to 0$
und $C\rho^{2} \to 0$, also $\Phi(\rho) \to A$ (endlich, von unten, da
$f(\rho) < 0$ f\"ur $\rho < \rho_{\mathrm{Pl}}$).
F\"ur $\rho \to +\infty$: $C\rho^{2}$ dominiert und $\Phi \to +\infty$.
Da $\Phi$ stetig auf $(0,\infty)$ ist, gegen~$A$ f\"ur $\rho \to 0^+$
und gegen~$+\infty$ f\"ur $\rho \to +\infty$ strebt, nimmt $\Phi$ sein
Infimum bei einem $\rho_{\ast} > 0$ an (das Infimum ist endlich und
echt kleiner als~$A$, also ein Minimum).

\medskip\noindent
\textbf{Schritt~2: Eindeutigkeit.}\;
Die zweite Ableitung ist
\begin{equation}
  \frac{\dd^{2}\Phi}{\dd\rho^{2}}
  \;=\; \int_{\LIR}^{\LUV}
  \left[\,
    \frac{T_{\mathrm{dS}}}{\rho}
    \;+\; \frac{16\pi G}{k^{2}}
  \,\right]
  \frac{4\pi k^{2}}{(2\pi)^{3}}\;\dd k
  \;=\; \frac{B}{\rho} + 2C
  \;>\; 0
  \quad\text{f\"ur alle } \rho > 0,
\end{equation}
da $B, C > 0$.  Strikte Konvexit\"at impliziert Eindeutigkeit des Minimums.

\medskip\noindent
\textbf{Schritt~3: Skalierung.}\;
Aus $\dd\Phi/\dd\rho = 0$ ergibt sich die Stationarit\"atsbedingung
\begin{equation}\label{eq:stationarity}
  B\,\bigl[1 + \ln(\rho_{\ast}/\rho_{\mathrm{Pl}})\bigr]
  \;+\; 2C\,\rho_{\ast} \;=\; 0.
\end{equation}
Da $\rho_{\ast} \ll \rho_{\mathrm{Pl}}$, ist der Logarithmus
$\ln(\rho_{\ast}/\rho_{\mathrm{Pl}})$ gross und negativ, sodass der
erste Term von $-B\,\ln(\rho_{\mathrm{Pl}}/\rho_{\ast})$ dominiert wird
und das Gleichgewicht lautet
\begin{equation}\label{eq:balance}
  \rho_{\ast} \;\approx\;
  \frac{B}{2C}\;\Bigl[\ln\!\left(\frac{\rho_{\mathrm{Pl}}}{\rho_{\ast}}\right)
  - 1\Bigr].
\end{equation}
Eine direkte Auswertung der Modenintegrale $B$ und~$C$ involviert das
Verh\"altnis $I_{1}/I_{2}$, wobei $I_{1} = \int \dd\mu(k) \sim \LUV^{3}$
und $I_{2} = \int k^{-2}\,\dd\mu(k) \sim \LUV$, sodass
$B/C \sim T_{\mathrm{dS}}\,\LUV^{2}/(8\pi G)
= H_0\,M_{\mathrm{Pl}}^{2}/(16\pi^{2}\,G)$.
In nat\"urlichen Einheiten ($G = M_{\mathrm{Pl}}^{-2}$) ergibt sich
$B/C \sim H_0\,M_{\mathrm{Pl}}^{4}/(16\pi^{2})$.

Die rohe Modensummen-Skalierung liefert $\rho_{\ast} \sim
H_0\,M_{\mathrm{Pl}}^{4}\,\ln(M_{\mathrm{Pl}}/H_0)$, was viel zu gross
ist. Dies spiegelt die bekannte UV-Sensitivit\"at von
Vakuumenergie-Berechnungen wider: Die Modensumme wird von UV-Moden
dominiert und liefert nicht direkt den beobachteten Wert. Wie in
Abschnitt~\ref{sec:renorm} diskutiert, erh\"alt man die physikalisch
relevante Aussage nach Renormierungsgruppen-Matching, wobei die
UV-divergenten Anteile in Gegenterme absorbiert werden und nur das
\emph{infrarote Laufen} \"uberlebt. Im renormierten Rahmenwerk entsteht
der logarithmische Faktor $\ln(M_{\mathrm{Pl}}/H)$ aus dem Laufen
zwischen $\mu = M_{\mathrm{Pl}}$ und $\mu = H_{0}$, was
\begin{equation}
  \rho_{\ast}^{\mathrm{ren}} \;\sim\;
  \frac{H_0^{2}}{8\pi G}\;
  \frac{1}{\ln(M_{\mathrm{Pl}}/H_0)}
\end{equation}
ergibt, was nach Einsetzen in~\eqref{eq:Leff} die
Schranke~\eqref{eq:bound} liefert. Die Herleitung dieser Skalierung
aus renormierten Gr\"o\ss en wird in Abschnitt~\ref{sec:renorm}
pr\"asentiert; die obige Skizze etabliert lediglich die strukturellen
Eigenschaften (Existenz, Eindeutigkeit, Wettbewerb zwischen
entropischen und gravitativen Termen).

\medskip\noindent
\emph{Ein rigoroser Beweis erfordert Kontrolle \"uber die Modensumme
\"uber alle Teilchenspezies des Standardmodells sowie eine
ordnungsgem\"a\ss e Renormierung der UV-Divergenzen. Dies wird in der
vollst\"andigen Version behandelt.}
\end{proof}

\begin{corollary}[Numerische Absch\"atzung]\label{cor:numerical}
Mit $H_0 \approx 67.4\;\mathrm{km\,s^{-1}\,Mpc^{-1}}
\approx 2.2 \times 10^{-18}\;\mathrm{s^{-1}}$ und
$\ln(M_{\mathrm{Pl}}/H_0) \approx 140$ erhalten wir
\begin{equation}
  \Leff \;\approx\;
  \frac{H_0^{2}}{c^{2}} \times \frac{1}{140}
  \;\approx\; 3.8 \times 10^{-53}\;\mathrm{m}^{-2},
\end{equation}
was innerhalb einer Gr\"o\ss enordnung des Planck-2018-Wertes
$\Lambda_{\mathrm{obs}} = (1.11 \pm 0.02) \times 10^{-52}\;
\mathrm{m}^{-2}$~\cite{Planck2018} liegt.
\end{corollary}

% ============================================================
\subsection{Verbindung zur Skalar-Tensor-Gravitation}\label{sec:scalar-tensor}\label{sec:saturation}

Der ph\"anomenologische R\"uckwirkungsterm~\eqref{eq:Vgrav} l\"asst
eine nat\"urliche Einbettung in eine Skalar-Tensor-Theorie der
Gravitation zu. Wir zeigen, dass das
Freie-Energie-Funktional~$\Phi(\rho)$ als effektive Beschreibung
eines gravitativen Lagrangians mit einem Skalarfeld~$\phi$ und einer
h\"oheren Kr\"ummungskorrektur entsteht:
\begin{equation}\label{eq:Lagrangian}
  \mathcal{L} \;=\;
  \frac{R + \gamma\,R^{2}}{16\pi G}
  \;+\; \frac{1}{2}\,(\partial\phi)^{2}
  \;-\; V_{\mathrm{PT}}(\phi)
  \;+\; \mathcal{L}_{\mathrm{m}},
\end{equation}
wobei~$\gamma > 0$ die Dimension $[\mathrm{Energie}]^{-2}$ hat (in nat\"urlichen Einheiten, $[G] = [\mathrm{Energie}]^{-2}$, sodass $[\gamma R^2/(16\pi G)] = [\mathrm{Energie}]^4$ wie gefordert) und das
Skalarpotential die P\"oschl--Teller-Form annimmt:
\begin{equation}\label{eq:VPT}
  V_{\mathrm{PT}}(\phi) \;=\;
  \frac{V_{0}}{\cosh^{2}(\phi/\phi_{0})}.
\end{equation}
Diese Wahl ist nicht willk\"urlich. Unter allen exakt l\"osbaren
quantenmechanischen Potentialen ist das P\"oschl--Teller-Potential das
\emph{einzige}, das ein Tangens-hyperbolicus-S\"attigungsprofil f\"ur
die Skalarfeld-Energiedichte erzeugt (siehe
Satz~\ref{prop:saturation} unten). Es ist zudem das Potential, das im
konformen Sektor der Starobinsky-$R^{2}$-Inflation~\cite{Starobinsky1980}
nach der Weyl-Reskalierung
$g_{\mu\nu} \to e^{-\sqrt{16\pi G/3}\;\phi}\,g_{\mu\nu}$ auftritt.

\medskip
Auf einem r\"aumlich flachen FLRW-Hintergrund mit Skalenfaktor~$a(t)$
lautet die Klein--Gordon-Gleichung f\"ur~$\phi$
\begin{equation}\label{eq:KG}
  \ddot{\phi} + 3H\dot{\phi}
  \;=\; -\frac{\dd V_{\mathrm{PT}}}{\dd\phi}
  \;=\; \frac{2V_{0}}{\phi_{0}}\;
  \frac{\tanh(\phi/\phi_{0})}{\cosh^{2}(\phi/\phi_{0})}.
\end{equation}
Definiert man den dimensionslosen Dunkle-Energie-Dichteparameter
$\Omega_{\Phi}(a) \equiv \rho_{\phi}(a)/\rho_{\mathrm{crit}}(a)$,
wobei $\rho_{\phi} = \tfrac{1}{2}\dot{\phi}^{2} + V_{\mathrm{PT}}(\phi)$,
so reduziert sich die Sp\"atzeit-Attraktordynamik von~\eqref{eq:KG}
auf die \emph{S\"attigungs-ODE}:
\begin{equation}\label{eq:saturation-ODE}
  \frac{\dd\Omega_{\Phi}}{\dd a}
  \;=\; \kappa\,
  \left[\,
    1 - \left(\frac{\Omega_{\Phi}}{\Phi_{0}}\right)^{\!2}
  \,\right],
\end{equation}
wobei $\kappa > 0$ ein Ratenparameter ist, der durch~$V_{0}$
und~$\phi_{0}$ bestimmt wird, und $\Phi_{0} = V_{0}/\rho_{\mathrm{crit},0}$
das asymptotische S\"attigungsniveau darstellt.
Gleichung~\eqref{eq:saturation-ODE} hat die exakte L\"osung
\begin{equation}\label{eq:tanh-solution}
  \Omega_{\Phi}(a) \;=\;
  \Phi_{0}\;\tanh\!\bigl[\kappa\,(a - a_{\mathrm{trans}})\bigr],
\end{equation}
wobei $a_{\mathrm{trans}}$ der Skalenfaktor ist, bei dem
$\Omega_{\Phi} = 0$ (die Epoche der Materie--Dunkle-Energie-Gleichheit).

\begin{proposition}[Stabiler de-Sitter-Endzustand]\label{prop:saturation}
Die S\"attigungs-ODE~\eqref{eq:saturation-ODE} hat genau zwei
Fixpunkte: $\Omega_{\Phi} = \pm\Phi_{0}$. Der Fixpunkt
$\Omega_{\Phi} = +\Phi_{0}$ ist ein globaler Attraktor f\"ur alle
Anfangsbedingungen $\Omega_{\Phi}(a_{i}) \in (-\Phi_{0},\,\Phi_{0})$.
Insbesondere gilt $\Omega_{\Phi}(a) \to \Phi_{0}$ f\"ur $a \to \infty$,
was einer asymptotischen de-Sitter-Phase mit
$H^{2} \to H_{0}^{2}\,\Phi_{0}$ entspricht. Es tritt keine
Big-Rip-Singularit\"at auf.
\end{proposition}

\begin{proof}
Die rechte Seite von~\eqref{eq:saturation-ODE} verschwindet bei
$\Omega_{\Phi} = \pm\Phi_{0}$ und ist positiv f\"ur
$|\Omega_{\Phi}| < \Phi_{0}$. Linearisierung um
$\Omega_{\Phi} = \Phi_{0}$: Setze $\delta = \Phi_{0} - \Omega_{\Phi}$,
dann gilt $\dd\delta/\dd a \approx -2\kappa\,\delta/\Phi_{0} < 0$,
was exponentielle Stabilit\"at best\"atigt. Die
L\"osung~\eqref{eq:tanh-solution} ist f\"ur alle~$a$ durch
$|\Omega_{\Phi}| \leq \Phi_{0}$ beschr\"ankt, sodass die Energiedichte
endlich bleibt und keine zuk\"unftige Singularit\"at auftritt.
\end{proof}

Die modifizierte Friedmann-Gleichung unter Einbeziehung der
Skalarfelddynamik lautet
\begin{equation}\label{eq:modified-Friedmann}
  H^{2}(a) \;=\; H_{0}^{2}\,
  \bigl[\,\Omega_{m}\,a^{-3} \;+\; \Omega_{\Phi}(a)\,\bigr],
\end{equation}
wobei $\Omega_{m} \approx 0.315$ und $\Omega_{\Phi}(a)$ durch~\eqref{eq:tanh-solution} gegeben ist. Dies ersetzt die kosmologische
Konstante $\Omega_{\Lambda} = \mathrm{const.}$ durch eine dynamische
Gr\"o\ss e, die glatt von $\Omega_{\Phi} \approx 0$ in der
materiedominierten \"Ara zu $\Omega_{\Phi} \to \Phi_{0} \approx 0.685$
zu sp\"aten Zeiten \"ubergeht.

\medskip
Die Verbindung zum Freie-Energie-Funktional ist nun transparent: Die
thermodynamische Schranke des Theorems~\ref{thm:main} bestimmt den
\emph{asymptotischen} Wert $\Phi_{0} \sim H_{0}^{2}/(Gc^{2})$, w\"ahrend
der Skalar-Tensor-Lagrangian~\eqref{eq:Lagrangian} die \emph{Dynamik}
liefert, die das Universum zu diesem Gleichgewicht treibt. Der
R\"uckwirkungsterm~$V_{\mathrm{grav}}$ in~\eqref{eq:Vgrav} ist die
N\"aherung des effektiven Potentials f\"ur die vollst\"andige
$R^{2} + V_{\mathrm{PT}}$-Dynamik im Slow-Roll-Regime. Der folgende
Satz macht diesen Zusammenhang rigoros, indem $V_{\mathrm{grav}}$ aus
der $f(R)$-Spurgleichung hergeleitet wird.

\begin{proposition}[Herleitung von $V_{\mathrm{grav}}$ aus ersten Prinzipien]
\label{prop:Vgrav-from-fR}
F\"ur $f(R) = R + \gamma R^2$-Gravitation mit Skalaronmasse
$m_s^2 = 1/(6\gamma)$ liefert die Spur der modifizierten
Einsteingleichungen die Skalarongleichung
\begin{equation}\label{eq:trace-eqn}
  R + 6\gamma\,\Box R \;=\; -8\pi G\,T.
\end{equation}
Im quasistatischen Limes auf sub-Hubble-Skalen ($\partial_t^2 \ll k^2/a^2$)
induziert eine Vakuumenergiedichte-St\"orung~$\delta\rho$ bei mitbewegter
Wellenzahl~$k$ eine gravitative R\"uckwirkungs-Energiedichte
\begin{equation}\label{eq:Vgrav-derived}
  V_{\mathrm{grav}}^{f(R)}[\rho,\,k]
  \;=\;
  \frac{8\pi G\,\rho^2}{k^2}
  \;\left[1 + \frac{1}{3}\,
  \frac{k^{2}/a^{2}}{k^{2}/a^{2} + m_{s}^{2}}\right].
\end{equation}
Insbesondere gilt:
\begin{itemize}
  \item F\"ur $k \ll a\,m_s$ (super-Compton-Skalen):
    $V_{\mathrm{grav}}^{f(R)} \to 8\pi G\,\rho^2/k^2$, was
    Gl.~\eqref{eq:Vgrav} und die reine ART exakt reproduziert.
  \item F\"ur $k \gg a\,m_s$ (sub-Compton-Skalen):
    $V_{\mathrm{grav}}^{f(R)} \to (4/3)\cdot 8\pi G\,\rho^2/k^2$,
    verst\"arkt um den Faktor~$4/3$ durch die skalare f\"unfte Kraft.
\end{itemize}
Der Faktor $4/3$ im sub-Compton-Limes ist eine bekannte Konsequenz des
zus\"atzlichen skalaren Freiheitsgrades in der $f(R)$-Gravitation~\cite{HuSawicki2007}
und beeinflusst nicht die parametrische Skalierung des Minimums (er verschiebt
den numerischen Vorfaktor von~$\rho_\ast$ um $O(1)$). Der
ph\"anomenologische Ansatz~\eqref{eq:Vgrav} entspricht dem
ART-normalisierten ($k \ll am_s$) Limes; der vollst\"andige
$k$-abh\"angige Ausdruck~\eqref{eq:Vgrav-derived} sollte f\"ur
$k \gtrsim am_s$ verwendet werden.
\end{proposition}

\begin{proof}
Man nehme die Spur der $f(R)$-Feldgleichungen
$f_R G_{\mu\nu} + (g_{\mu\nu}\Box - \nabla_\mu\nabla_\nu)f_R
= 8\pi G\,T_{\mu\nu}$ mit $f_R = 1 + 2\gamma R$, $f_{RR} = 2\gamma$.
Die Spur ergibt~\eqref{eq:trace-eqn} (vgl.~Sotiriou \&
Faraoni~2010).

Im Fourierraum auf einem r\"aumlich flachen FLRW-Hintergrund, linearisiert
um $R = R_0$, induziert eine St\"orung $\delta\rho$ bei mitbewegter
Wellenzahl~$k$
\[
  \bigl(k^2/a^2 + m_s^2\bigr)\,\delta R
  \;=\; \frac{8\pi G}{3}\;\delta\rho,
\]
wobei $m_s^2 = 1/(6\gamma)$ das Quadrat der Skalaronmasse ist. Die
modifizierte Poissongleichung im quasistatischen, sub-Horizont-Regime
liefert~\cite{HuSawicki2007}
\[
  k^2\,\Phi_N \;=\; -4\pi G\,a^2\,\rho\,\delta
  \;\left[1 + \frac{1}{3}\,\frac{k^2/a^2}{k^2/a^2 + m_s^2}\right].
\]
Identifikation der gravitativen Selbstenergie-Dichte auf der Skala~$k$
als $V_{\mathrm{grav}}^{f(R)} = 8\pi G_{\mathrm{eff}}(k)\,\rho^2/k^2$
mit
\[
  G_{\mathrm{eff}}(k) \;=\; G\,\left[1 + \frac{1}{3}\,
  \frac{k^2/a^2}{k^2/a^2 + m_s^2}\right]
\]
liefert~\eqref{eq:Vgrav-derived} direkt.
F\"ur $k \ll a\,m_s$: $G_{\mathrm{eff}} \to G$ und
$V_{\mathrm{grav}}^{f(R)} \to 8\pi G\rho^2/k^2$, in
\"Ubereinstimmung mit~\eqref{eq:Vgrav}.
F\"ur $k \gg a\,m_s$: $G_{\mathrm{eff}} \to (4/3)\,G$, was eine
$4/3$-Verst\"arkung gegen\"uber dem ART-Wert ergibt.
\end{proof}

\begin{remark}[Schlie\ss ung der R\"uckwirkungsl\"ucke]
Satz~\ref{prop:Vgrav-from-fR} liefert die Herleitung
von~$V_{\mathrm{grav}}$ aus ersten Prinzipien, die als zentrales
offenes Problem~(L1) in der Review-Analyse identifiziert wurde.
Zusammen mit Lemma~\ref{lem:minimality} (welches die Eindeutigkeit aus
dimensionalen Axiomen begr\"undet) ist $V_{\mathrm{grav}}$ nun durch zwei
unabh\"angige Argumente fundiert: eine \emph{konstruktive} Herleitung aus
der $f(R)$-Lagrangedichte und ein \emph{Eindeutigkeitsbeweis} aus der
Dimensionsanalyse. Die $f(R)$-Herleitung reproduziert
$V_{\mathrm{grav}} = 8\pi G\rho^2/k^2$ exakt im super-Compton-Limes
(ART) und sagt eine $4/3$-Verst\"arkung auf sub-Compton-Skalen voraus.
Da das Freie-Energie-Minimum \"uber alle Moden von $\LIR$ bis $\LUV$
integriert, betrifft der $4/3$-Faktor nur den sub-Compton-Anteil des
Integrals und \"andert den numerischen Koeffizienten von~$\rho_\ast$
um $O(1)$, ohne die parametrische Skalierung~\eqref{eq:bound} zu
\"andern.
\end{remark}

% ============================================================
\subsection{Selbstkonsistente Dynamik und G\"ultigkeit des S\"attigungsansatzes}
\label{sec:self-consistent}

Die S\"attigungs-ODE~\eqref{eq:saturation-ODE} ist eine effektive
Reduktion des vollst\"andig gekoppelten Systems aus Friedmann-,
Klein--Gordon- und Materie-Kontinuit\"atsgleichungen. Wir stellen nun
das geschlossene System vor und etablieren Bedingungen, unter denen
die Reduktion kontrolliert ist.

\medskip
\noindent\textbf{Autonome Formulierung.}\;
Definiere $N := \ln a$ als die Anzahl der $e$-Faltungen und f\"uhre
dimensionslose Phasenraumvariablen ein~(vgl.~\cite{CopelandLiddleWands1998}):
\begin{equation}\label{eq:xy-def}
  x \;:=\; \frac{\dot\phi}{\sqrt{6}\,H\,M_{\mathrm{Pl}}},
  \qquad
  y \;:=\; \frac{\sqrt{V}}{\sqrt{3}\,H\,M_{\mathrm{Pl}}}.
\end{equation}
Der Dunkle-Energie-Dichteparameter und der Zustandsgleichungsparameter
sind $\Omega_\phi = x^2 + y^2$ und $w_\phi = (x^2 - y^2)/(x^2 + y^2)$.
F\"ur das P\"oschl--Teller-Potential~\eqref{eq:VPT} lautet der
Steigungsparameter
\begin{equation}\label{eq:lambda-PT}
  \lambda(\phi) \;:=\;
  -\frac{M_{\mathrm{Pl}}\,V'(\phi)}{V(\phi)}
  \;=\; \frac{2\,M_{\mathrm{Pl}}}{\phi_0}\,
  \tanh\!\left(\frac{\phi}{\phi_0}\right).
\end{equation}
In diesen Variablen wird das Friedmann--Klein--Gordon-System mit
druckloser Materie ($w_m = 0$) zum \emph{autonomen} System
\begin{align}
  \frac{dx}{dN} &\;=\;
    -3x + \frac{\sqrt{6}}{2}\,\lambda(\phi)\,y^2
    + \frac{3}{2}\,x\,(1 - \Omega_\phi),
    \label{eq:dx}\\[4pt]
  \frac{dy}{dN} &\;=\;
    -\frac{\sqrt{6}}{2}\,\lambda(\phi)\,x\,y
    + \frac{3}{2}\,y\,(1 - \Omega_\phi),
    \label{eq:dy}\\[4pt]
  \frac{d\phi}{dN} &\;=\; \sqrt{6}\,x\,M_{\mathrm{Pl}}.
    \label{eq:dphi}
\end{align}
Der Hubble-Parameter ergibt sich aus der Friedmann-Bedingung
$H^2 = \rho_m/(3M_{\mathrm{Pl}}^2\,(1 - \Omega_\phi))$.

\medskip
\noindent\textbf{Reduktion auf die S\"attigungs-ODE.}\;
Die exakte Entwicklungsgleichung f\"ur $\Omega_\phi$ folgt aus
\eqref{eq:dx}--\eqref{eq:dy}:
\begin{equation}\label{eq:dOmega-exact}
  \frac{d\Omega_\phi}{dN}
  \;=\; 3\,(1 + w_\phi)\,\Omega_\phi\,(1 - \Omega_\phi).
\end{equation}
Dies ist logistisch in $\Omega_\phi$, wann immer $1 + w_\phi$
n\"aherungsweise konstant ist. Im \emph{Thawing}-Regime---in dem $\phi$
nahe dem Maximum von~$V_{\mathrm{PT}}$ bei $\phi \approx 0$ mit
$\dot\phi \approx 0$ startet, sodass $x^2 \ll y^2$ und
$w_\phi \approx -1$---variiert die Gr\"o\ss e
$1 + w_\phi \approx 2x^2/(x^2 + y^2)$ langsam. Mit der Identifikation
$\kappa = 3(1 + w_\phi)/(2\Phi_0)$ und $\Phi_0 = \Omega_\phi^{(\infty)}$
reduziert sich Gleichung~\eqref{eq:dOmega-exact} auf die
S\"attigungs-ODE~\eqref{eq:saturation-ODE} mit $dN = d(\ln a) = da/a$.

\medskip
\noindent\textbf{Fehlerkontrolle.}\;
Schreibe die exakte Dynamik als
\begin{equation}\label{eq:Omega-rest}
  \frac{d\Omega_\phi}{da}
  \;=\; \kappa\!\left[1 - \left(\frac{\Omega_\phi}{\Phi_0}\right)^{\!2}\right]
  \;+\; R(a),
\end{equation}
wobei das Residuum $R(a)$ die Variation von $w_\phi(a)$ und
$\lambda(\phi(a))$ gegen\"uber ihren Bezugswerten erfasst. Nach der
Gr\"onwall-Ungleichung erf\"ullt die Differenz zwischen der exakten
L\"osung $\Omega_\phi^{\mathrm{full}}(a)$ und der
tanh-N\"aherung~\eqref{eq:tanh-solution}
\begin{equation}\label{eq:Gronwall-error}
  \bigl|\Omega_\phi^{\mathrm{full}}(a)
  - \Omega_\phi^{\mathrm{tanh}}(a)\bigr|
  \;\le\;
  \int_{a_i}^{a} |R(u)|\,
  \exp\!\left(\int_u^a L(v)\,dv\right) du,
\end{equation}
wobei $L(a) = 2\kappa\,|\Omega_\phi|/\Phi_0^2$ die Lipschitz-Konstante
der rechten Seite von~\eqref{eq:saturation-ODE} ist. F\"ur das
P\"oschl--Teller-Potential ist das Residuum beschr\"ankt durch
\begin{equation}
  |R(a)| \;\le\;
  C\,\sup_{u \in [a_i,\,a]}|w_\phi'(u)|
  + C'\,\sup_{u \in [a_i,\,a]}|\lambda'(\phi(u))|,
\end{equation}
was beides im Thawing-Regime, wo $|\lambda(\phi)| \ll 1$, d.\,h.\
$\phi_0 \gg M_{\mathrm{Pl}}$, klein ist.

\medskip
\noindent\textbf{Parameteridentifizierbarkeit.}\;
Das Modell hat drei freie Parameter $(V_0, \phi_0, \gamma)$.
Mit Thawing-Anfangsbedingungen ($\dot\phi \approx 0$, $\phi \approx 0$
zu fr\"uhen Zeiten) bestimmen die Beobachtungsbeschr\"ankungen
$\Omega_\Lambda = 0.685$, $z_{\mathrm{acc}} \approx 0.7$ und
$H_0 = 67.4\;\mathrm{km\,s^{-1}\,Mpc^{-1}}$ diese Parameter wie folgt:
\begin{itemize}
  \item $V_0$ wird durch die Forderung
        $\Omega_\phi(a=1) = 0.685$ fixiert, was
        $V_0 \approx 3M_{\mathrm{Pl}}^2 H_0^2\,\Omega_\Lambda$ ergibt.
  \item $\phi_0$ kontrolliert die \"Ubergangsbreite \"uber $\lambda(\phi)$
        und wird durch $z_{\mathrm{acc}} \approx 0.7$ fixiert.
  \item $\gamma$ geht nur \"uber die Skalaronmasse
        $m_s^2 = 1/(6\gamma)$ in die Hintergrunddynamik ein. F\"ur
        Sp\"atzeit-Dunkle-Energie ist das Skalaron schwer und entkoppelt;
        $\gamma$ wird nat\"urlicherweise eher durch Fr\"uhuniversums-Daten
        (CMB, Inflation) oder den
        Gravitations-Slip-Parameter~\eqref{eq:grav-slip} eingeschr\"ankt.
\end{itemize}
Ohne Fixierung der Anfangsbedingungen treten Entartungen zwischen
$\phi_0$ und $\gamma$ auf; die Thawing-Annahme bricht diese Entartung
und macht $(V_0, \phi_0)$ praktisch allein aus Sp\"atzeit-Daten
identifizierbar.

\medskip
\noindent\textbf{Stabilit\"at des de-Sitter-Endzustands.}\;
Der de-Sitter-Fixpunkt entspricht $x = 0$, $y = 1$,
$\Omega_\phi = 1$, $\lambda = 0$ (d.\,h.\ $\phi = 0$, das Maximum
von~$V_{\mathrm{PT}}$). Linearisierung der Klein--Gordon-Gleichung um
$\phi = 0$:
\begin{equation}
  \ddot\phi + 3H\dot\phi + V''(0)\,\phi \;=\; 0,
  \qquad
  V''(0) \;=\; -\frac{2V_0}{\phi_0^2} \;<\; 0.
\end{equation}
Die negative Kr\"ummung $V''(0) < 0$ ist ein tachyonischer Term, aber
die Hubble-Reibung liefert \"Uberd\"ampfung, wenn
\begin{equation}\label{eq:stability-condition}
  |V''(0)| \;\le\; \frac{9}{4}\,H^2
  \qquad\Longleftrightarrow\qquad
  \phi_0 \;\ge\; \sqrt{\frac{8}{3}}\;M_{\mathrm{Pl}}
  \;\approx\; 1.63\,M_{\mathrm{Pl}}.
\end{equation}
Unter dieser Bedingung---die einen super-Planckschen Feldbereich
erfordert, konsistent mit Gro\ss feld-Inflationsmodellen---werden kleine
St\"orungen um $\phi = 0$ exponentiell ged\"ampft und der
de-Sitter-Attraktor von Satz~\ref{prop:saturation} ist im
vollst\"andigen Klein--Gordon--Friedmann-System stabil. Eine
vollst\"andige Stabilit\"atsanalyse einschlie\ss lich
Metrikst\"orungen erfordert den auf den Skalar-Tensor-Sektor
erweiterten Mukhanov--Sasaki-Formalismus und wird auf zuk\"unftige
Arbeiten verschoben.

% ============================================================
\section{Vergleich mit Beobachtungen}\label{sec:observations}

\subsection{Konsistenz mit Planck 2018}

Die Messung des Dunkle-Energie-Dichteparameters durch den
Planck-Satelliten ergibt
$\Omega_{\Lambda} = 0.6847 \pm 0.0073$~\cite{Planck2018}.
Kombiniert mit $H_0 = 67.36 \pm 0.54\;\mathrm{km\,s^{-1}\,Mpc^{-1}}$
folgt daraus
\begin{equation}
  \Lambda_{\mathrm{obs}} \;=\; 3\,\Omega_{\Lambda}\,
  \frac{H_0^{2}}{c^{2}}
  \;\approx\; 1.11 \times 10^{-52}\;\mathrm{m}^{-2}.
\end{equation}
Unsere Absch\"atzung aus Korollar~\ref{cor:numerical} liefert
$\Leff \approx 3.8 \times 10^{-53}\;\mathrm{m}^{-2}$, einen Faktor
$\sim 3$ unter dem beobachteten Wert. Angesichts der Tatsache, dass die
Entwurfsherleitung numerische Vorfaktoren und speziesabh\"angige
Beitr\"age vernachl\"assigt, ist dieses \"Ubereinstimmungsniveau
ermutigend.

\subsection{Das Koinzidenzproblem}

Ein anhaltendes R\"atsel in der Dunkle-Energie-Kosmologie ist, warum
$\Omega_{\Lambda} \sim \Omega_{m}$ heute gilt -- das
\emph{Koinzidenzproblem}. In unserem Rahmenwerk findet dies eine
nat\"urliche Erkl\"arung: Der Infrarot-Cutoff~$\LIR = H_0$ entwickelt
sich mit der Expansion, sodass $\Leff(t)$ dem Wert~$H(t)^{2}$ folgt.
W\"ahrend der Materiedominanz gilt $H^{2} \propto \rho_{m}$, was
automatisch $\rho_{\mathrm{vac}} \sim \rho_{m}$ in der Epoche
sicherstellt, in der der \"Ubergang von Verz\"ogerung zu Beschleunigung
stattfindet.

Die S\"attigungsdynamik aus Abschnitt~\ref{sec:scalar-tensor} macht
dieses Argument pr\"azise. Aus der modifizierten
Friedmann-Gleichung~\eqref{eq:modified-Friedmann} ergibt sich der
Verz\"ogerungsparameter
\begin{equation}
  q(a) \;=\; \frac{\Omega_{m}\,a^{-3}}{2\,[\Omega_{m}\,a^{-3}
  + \Omega_{\Phi}(a)]}
  \;-\; \frac{\Omega_{\Phi}(a) + a\,\Omega_{\Phi}'(a)/3}
       {\Omega_{m}\,a^{-3} + \Omega_{\Phi}(a)}.
\end{equation}
Setzt man $q(a_{\mathrm{acc}}) = 0$ mit der
tanh-L\"osung~\eqref{eq:tanh-solution} und den illustrativen Parametern
$\kappa \approx 1{,}5$, $a_{\mathrm{trans}} \approx 0{,}50$, so erh\"alt man
$a_{\mathrm{acc}} \approx 0.59$, entsprechend
$z_{\mathrm{acc}} \approx 0.70$, in hervorragender \"Ubereinstimmung
mit dem beobachteten Wert.

\subsection{Effektive Zustandsgleichung}\label{sec:weff}

Die tanh-Dynamik aus Abschnitt~\ref{sec:scalar-tensor} bestimmt den
effektiven Zustandsgleichungsparameter f\"ur die Skalarfeld-Dunkle-Energie.
Mit der Definition
\begin{equation}\label{eq:weff}
  w_{\mathrm{eff}}(a)
  \;=\; -1 \;-\; \frac{1}{3}\;
  \frac{\dd\ln\Omega_{\Phi}}{\dd\ln a},
\end{equation}
und Einsetzen der tanh-L\"osung~\eqref{eq:tanh-solution}:
\begin{equation}
  \frac{\dd\ln\Omega_{\Phi}}{\dd\ln a}
  \;=\; \frac{a\,\kappa}
       {\Phi_{0}\,\tanh[\kappa(a - a_{\mathrm{trans}})]}\;
  \Bigl[1 - \tanh^{2}\!\bigl[\kappa(a - a_{\mathrm{trans}})\bigr]
  \Bigr].
\end{equation}
Auswertung bei $a = 1$ (d.\,h.\ $z = 0$) mit illustrativen Parametern
$\kappa \approx 1{,}5$--$2{,}0$, $a_{\mathrm{trans}} \approx 0{,}50$,
$\Phi_{0} \approx 0{,}685$:
\begin{equation}\label{eq:w0-prediction}
  w_{\mathrm{eff}}(z=0) \;\approx\; -1{,}4 \;\text{bis}\; -1{,}5.
\end{equation}
Dies platziert das Modell im \emph{Phantomregime}
($w_{\mathrm{eff}} < -1$), was auf den ersten Blick besorgniserregend
erscheinen mag.
(Der pr\"azise Wert h\"angt von der Parameteranpassung ab; ein
dedizierter numerischer Fit an Supernova- und BAO-Daten ist
erforderlich, um $\kappa$ und $a_{\mathrm{trans}}$ auf besser als
$\sim 30\%$ zu fixieren.)

\medskip\noindent
\textbf{Wichtige Einschr\"ankung.}\;
Die Gr\"o\ss e~$w_{\mathrm{eff}}$ gem\"a\ss~\eqref{eq:weff} ist ein
\emph{effektiver} Zustandsgleichungsparameter, der den Effekt der
zeitlichen Variation von~$\Omega_{\Phi}$ einschlie\ss t; er unterscheidet
sich von der intrinsischen Zustandsgleichung der Dunklen Energie
$w_{\mathrm{DE}} = p_{\phi}/\rho_{\phi}$. Im Thawing-Regime gilt
$w_{\mathrm{DE}} \approx -1 + 2x^{2}/(x^{2}+y^{2}) \gtrsim -1$,
das Skalarfeld selbst ist also \emph{kein} Phantom.
Der effektive Phantomwert $w_{\mathrm{eff}} < -1$ entsteht, weil
$\Omega_{\Phi}$ w\"achst (der Dunkle-Energie-Anteil nimmt auf Kosten
der Materie zu), nicht wegen eines geistartigen kinetischen Terms.

Die DESI-DR2-Analyse~\cite{DESI2025} berichtet
$w_{0} = -0.42 \pm 0.21$ und $w_{a} = -1.79 \pm 0.67$ in der
$w_{0}w_{a}$CDM-Parametrisierung. Ein direkter Vergleich mit unserem
$w_{\mathrm{eff}}(z=0) \approx -1{,}4$ bis $-1{,}5$ ist nicht unmittelbar m\"oglich,
da das DESI-$w_{0}$ einer anderen Parametrisierung entspricht.
Ein korrekter Vergleich erfordert die Projektion unserer tanh-Dynamik
auf die $(w_{0}, w_{a})$-Ebene, was einer gesonderten numerischen
Studie vorbehalten bleibt.

Entscheidend ist, dass das
Phantomverhalten in unserem Modell \emph{nicht} pathologisch ist:

\begin{itemize}
  \item Die Dunkle-Energie-Dichte $\Omega_{\Phi}(a)$ ist f\"ur alle~$a$
        nach oben durch $\Phi_{0}$ beschr\"ankt
        (Satz~\ref{prop:saturation}).
  \item Die Phantomphase ist vor\"ubergehend: F\"ur $a \to \infty$ gilt
        $\Omega_{\Phi} \to \Phi_{0}$ und
        $w_{\mathrm{eff}} \to -1$ von unten.
  \item Es tritt keine Big-Rip-Singularit\"at auf, im Gegensatz zu
        Phantom-Skalarfeldmodellen mit unbeschr\"ankter kinetischer
        Energie~\cite{Carroll2001}.
\end{itemize}
Der S\"attigungsmechanismus bietet somit eine nat\"urliche Aufl\"osung
des Problems der ,,Phantom-Divide-\"Uberquerung``, das viele dynamische
Dunkle-Energie-Modelle plagt.

\begin{remark}[Numerischer Vergleich mit DESI DR2]\label{rem:desi-comparison}
Ein Gitterscan \"uber den Parameterraum $(\kappa, a_{\mathrm{trans}}) \in [0{,}5, 5] \times [0{,}3, 0{,}8]$ wurde durchgef\"uhrt, wobei die intrinsische Zustandsgleichung der Dunklen Energie $w_{\mathrm{DE}}(z)$ (nicht $w_{\mathrm{eff}}$) aus der Friedmann-Kontinuit\"atsgleichung berechnet und an die CPL-Parametrisierung $w(z) = w_0 + w_a z/(1+z)$ angepasst wurde.

\smallskip\noindent\textbf{Bester Fit an DESI DR2 (2025):} $\kappa = 2{,}63$, $a_{\mathrm{trans}} = 0{,}326$ (entsprechend $z_{\mathrm{trans}} \approx 2{,}07$), mit $w_0^{\mathrm{DE}} = -0{,}470$ und $w_a^{\mathrm{DE}} = -2{,}00$, bei $\chi^2 = 0{,}15$ ($0{,}39\sigma$ vom DESI-Bestwert $w_0 = -0{,}42 \pm 0{,}21$, $w_a = -1{,}79 \pm 0{,}67$). Der kompatible Bereich ($<2\sigma$) umfasst $\kappa \in [1{,}2, 5{,}0]$ und $a_{\mathrm{trans}} \in [0{,}30, 0{,}33]$.

\smallskip\noindent\textbf{Schl\"usseleinsicht:} Der \"Ubergangsskalenfaktor muss $a_{\mathrm{trans}} \lesssim 0{,}33$ ($z_{\mathrm{trans}} \gtrsim 2$) erf\"ullen, damit die $\tanh$-Singularit\"at von $\Omega_\Phi$ au\ss erhalb des beobachtungskontrollierten Rotverschiebungsbereichs liegt. Mit den illustrativen Parametern $\kappa = 1{,}75$, $a_{\mathrm{trans}} = 0{,}5$ aus Abschnitt~\ref{sec:saturation} liegt die Singularit\"at bei $z = 1$, was den CPL-Vergleich pathologisch macht. Dies best\"atigt, dass die beobachtungsbezogene Tragf\"ahigkeit des Modells prim\"ar von $a_{\mathrm{trans}}$ abh\"angt, nicht von $\kappa$.

\smallskip\noindent Numerische Skripte: \texttt{compute\_w\_mapping.py}.
\end{remark}

\subsection{BBN-Schutz durch Spurkopplung}\label{sec:BBN}

Ein h\"aufiges Bedenken bei modifizierten Gravitationstheorien ist ihr
potenzieller Einfluss auf die Urknall-Nukleosynthese (BBN), die die
Expansionsrate bei $T \sim 1\;\mathrm{MeV}$ auf Sub-Prozent-Niveau
einschr\"ankt~\cite{SotiriouFaraoni2010}. Im
Skalar-Tensor-Lagrangian~\eqref{eq:Lagrangian} erzeugt die
$R^{2}$-Korrektur einen effektiven skalaren Freiheitsgrad (das
Skalaron), der an Materie \"uber die Spur des
Energie-Impuls-Tensors koppelt:
\begin{equation}\label{eq:trace}
  T \;=\; g^{\mu\nu}\,T_{\mu\nu}.
\end{equation}
F\"ur ein perfektes Fluid mit Zustandsgleichungsparameter~$w$ gilt f\"ur die
Spur $T = (\rho - 3p) = (1 - 3w)\,\rho$. W\"ahrend der
strahlungsdominierten Epoche ($w = 1/3$):
\begin{equation}\label{eq:trace-rad}
  T_{\mathrm{rad}} \;=\; (1 - 3 \times \tfrac{1}{3})\,\rho_{\mathrm{rad}}
  \;=\; 0.
\end{equation}
Das Verschwinden der Spur ist eine Konsequenz der konformen Symmetrie
massloser Strahlung. Da das $R^{2}$-Skalaron \"uber die
Spurgleichung~\eqref{eq:trace-eqn} an Materie koppelt, ist der
Quellterm f\"ur Skalaron-Anregungen proportional zu~$T$ und daher
w\"ahrend der Strahlungs\"ara stark unterdr\"uckt. Genauer impliziert
die Spurgleichung $R + 6\gamma\Box R = -8\pi GT$, dass die
Skalaron-Amplitude $\delta R \propto T/(k^2/a^2 + m_s^2)$ durch~$T$
getrieben wird; f\"ur $T_{\mathrm{rad}} = 0$ (ideale masselose
Strahlung) entkoppelt das Skalaron vollst\"andig. Residuale
Beitr\"age von massiven Spezies ($e^{\pm}$, $\nu$ nach Entkopplung)
und der Spuranomalie sind durch $m^2/T^2$ bzw.\ $\alpha_s^2$
unterdr\"uckt und treten w\"ahrend der BBN ($T \sim 1\;\mathrm{MeV}$)
nur auf Sub-Prozent-Niveau ein.
Diese strukturelle Unterdr\"uckung stellt sicher, dass die
Skalar-Tensor-Dynamik weder das Neutronen-Protonen-Verh\"altnis noch
die primordiale Heliumh\"aufigkeit wesentlich beeinflusst, ohne dass
eine Feinabstimmung von~$\gamma$ erforderlich ist. Der BBN-Schutz
ist eine Konsequenz f\"uhrender Ordnung der konformen Symmetrie und der
Spurkopplungsstruktur der Theorie.

\begin{remark}[Sonnensystem-Einschr\"ankungen]\label{sec:screening}
F\"ur eine Skalaronmasse $m_{s} \sim H_{0} \sim 10^{-33}\;\mathrm{eV}$
ist die Compton-Wellenl\"ange $\lambda_{C} \sim c/m_{s}$ von der
Gr\"o\ss enordnung des Hubble-Radius. Auf Sonnensystem-Skalen
($r \sim 1\;\mathrm{AU} \ll \lambda_{C}$) ist das Skalarfeld effektiv
masselos und vermittelt eine f\"unfte Kraft mit Kopplungsst\"arke
$\sim G/3$, die den PPN-Parameter $\gamma_{\mathrm{PPN}}$ auf
$\gamma_{\mathrm{PPN}} = 1/2$ modifiziert---stark ausgeschlossen
durch Cassini-Tracking
($|\gamma_{\mathrm{PPN}} - 1| < 2{,}3 \times 10^{-5}$).
Hu \& Sawicki~\cite{HuSawicki2007} zeigten, dass kosmologisch
viable $f(R)$-Modelle ein nichtlineares Potential erfordern, bei dem
$m_s \gg H_0$ in Hochdichte-Umgebungen gilt (Chameleon-Mechanismus).
Im vorliegenden Rahmenwerk platziert das
$\gamma \gg 10^{60}$-Regime (Abschnitt~\ref{sec:one-loop}) die
Skalaronmasse weit unter $H_{0}$, sodass der lineare
$R^2$-Lagrangian~\eqref{eq:Lagrangian} \emph{kein}
Chameleon-Screening bietet. Zwei m\"ogliche Aufl\"osungen bestehen:
(i)~Ersetzung des linearen $R^2$-Terms durch ein nichtlineares $f(R)$
vom Hu--Sawicki-Typ, was den UV-Sektor modifizieren, aber die
IR-Skalierung des Minimums bewahren w\"urde;
(ii)~Interpretation der $R^2$-Einbettung als
effektive Theorie, die nur auf kosmologischen Skalen g\"ultig ist,
wobei die Kurzstreckenphysik durch einen Vainshtein-Mechanismus aus
dem Skalar-Tensor-Sektor abgeschirmt wird.
Wir betrachten die Konstruktion eines viablen Screening-Mechanismus
als die zentrale offene Herausforderung f\"ur die
Skalar-Tensor-Einbettung und verschieben sie auf zuk\"unftige
Arbeiten.
\end{remark}

\subsection{Screening als $\Gamma$-konvergente Phasenseparation}\label{sec:screening-gamma}

Das Fifth-Force-Problem ist keine externe Randbedingung, sondern ein \emph{Grenzph\"anomen} eines nichtkonvexen thermodynamischen Funktionals. Screening tritt als Phasen\"ubergang im $\Gamma$-Grenzwert auf, nicht als postulierter Mechanismus.

\subsubsection*{Schritt 1: Von konvexer zu nichtkonvexer effektiver Energie}

Die mikroskopische freie Energie $U(\rho) = \rho\ln(\rho/\rho_{\mathrm{Pl}})$ aus den Abschnitten~\ref{sec:functional}--\ref{sec:selection} ist \emph{konvex}, was einen eindeutigen Minimierer $\Lambda_{\mathrm{eff}}$ sicherstellt, aber kein Screening erzeugt. Um die gravitative R\"uckreakti\-on einzubeziehen, gehen wir zu einer \emph{effektiven} Energie \"uber, die die Materie-Skalar-Kopplung enth\"alt:
\begin{equation}\label{eq:phi-eps-screening}
  \Phi_\varepsilon[\rho, \varphi] = \int \Bigl( U(\rho) + V_{\mathrm{eff}}(\varphi;\rho) + \varepsilon\,|\nabla\varphi|^2 \Bigr)\,dx,
\end{equation}
wobei $\varphi$ das Skalaronfeld ist und $V_{\mathrm{eff}}(\varphi;\rho) = V(\varphi) + \rho\,A(\varphi)$ ein dichteabh\"angiges effektives Potential ist. F\"ur geeignete $V$ und $A$ (z.\,B.\ Hu--Sawicki \cite{HuSawicki2007} oder Symmetron-Modelle) entwickelt $V_{\mathrm{eff}}$ bei hoher Umgebungsdichte $\rho_{\mathrm{env}}$ \emph{zwei Minima}: ein tiefes Minimum (schweres Skalaron, abgeschirmt) und ein flaches Minimum (leichtes Skalaron, nicht abgeschirmt). Die Nichtkonvexit\"at entsteht nicht in $U(\rho)$ selbst, sondern im \emph{gemeinsamen} Funktional $\Phi_\varepsilon[\rho,\varphi]$ durch die Materie-Skalar-Kopplung.

\subsubsection*{Schritt 2: $\Gamma$-Konvergenz und Phasenseparation}

F\"ur $\varepsilon \to 0$:
\[
  \Phi_\varepsilon \;\xrightarrow{\;\Gamma\;}\; \Phi_0,
\]
wobei $\Phi_0$ ein \emph{phasensepariertes} Funktional ist. Die Konsequenzen sind:
\begin{itemize}
\item \textbf{Niedrige Umgebungsdichte} (kosmologische Skalen): $\rho$ minimiert die flache Phase $\Rightarrow$ effektive dunkle Energie $\Lambda_{\mathrm{eff}} > 0$.
\item \textbf{Hohe Umgebungsdichte} (Sonnensystem): $\rho$ kollabiert in die steile Phase $\Rightarrow$ das Skalaron wird effektiv massiv, $m_{\mathrm{eff}} \gg H_0$.
\end{itemize}
Dies ist \emph{pr\"azise} der Cham\"aleon-Effekt, jedoch aus der Variationsstruktur hergeleitet statt als Mechanismus postuliert.

\subsubsection*{Schritt 3: Energie-Dissipations-Ungleichung}

F\"ur Minimierer $\rho_\varepsilon$ von $\Phi_\varepsilon$ liefert die EDP-Struktur:
\begin{equation}\label{eq:edp-screening}
  \varepsilon \int |\nabla\rho_\varepsilon|^2 \;\leq\; \Phi_\varepsilon[\rho_\varepsilon] - \inf\Phi_\varepsilon.
\end{equation}
Im Grenzwert:
\begin{itemize}
\item verschwindet die Gradientenenergie,
\item schrumpfen die \"Ubergangszonen,
\item wird die effektive Kopplung an Materie singul\"ar unterdr\"uckt.
\end{itemize}
\emph{Screening ist kein dynamischer Prozess, sondern ein thermodynamischer Phasen\"ubergang.}

\subsubsection*{Schritt 4: Beziehung zu Hu--Sawicki-Modellen}

Hu--Sawicki-$f(R)$-Theorien sind eine spezifische parametrische Darstellung des $\Gamma$-Grenzwerts von $\Phi_\varepsilon$. Der wesentliche Unterschied: im Standard-$f(R)$ ist Screening eingebaut; in der FST-Formulierung wird Screening \emph{erzwungen}, weil andernfalls $\Phi_0$ nicht minimiert wird. Das $\Lambda$-Selektionsprinzip der Abschnitte~\ref{sec:functional}--\ref{sec:selection} bleibt daher vollst\"andig erhalten.

\begin{remark}[Autonome $\Lambda$-Selektion]\label{rem:autonomous-lambda}
Die Abschnitte~\ref{sec:functional}--\ref{sec:selection} (thermodynamische $\Lambda$-Selektion) sind eigenst\"andig und unabh\"angig von der $f(R)$-Einbettung publikationsf\"ahig. $\Lambda$ ist ein thermodynamischer Ordnungsparameter; Gravitation ist ein effektiver Langwellen-Tr\"ager; Fifth-Force-Probleme entstehen nur durch falsche Grenzwert-Reihenfolge.
\end{remark}

\begin{remark}[Problemverbindende Struktur]\label{rem:cross-edp-de}
Dieser $\Gamma$/EDP-Mechanismus ist strukturell identisch mit:
\begin{itemize}
\item BV-Selektion durch integrierte Dissipation (Navier--Stokes, \cite{FST-NS2026}),
\item RG-Fluss-Kontraktion via subadditive Ergodentheorie (Yang--Mills, \cite{Geiger2026-YM}),
\item H\"ohens\"attigung als Ersatz algebraischer Kontrolle (BSD, \cite{Geiger2026-BSD}).
\end{itemize}
Alle teilen das Prinzip: \emph{lokale Dynamik $+$ Dissipation $\Rightarrow$ globale Selektion im Grenzwert}.
\end{remark}

\begin{remark}[Screening-Phasenseparation]
\label{rem:screening-numerics}
Das gemeinsame Funktional $\Phi_\varepsilon[\rho, \varphi]$ aus Abschnitt~\ref{sec:screening-gamma} wurde in 1D mit Hu--Sawicki-Potential $V(\varphi)$ und konformaler Kopplung $A(\varphi) = e^{\varphi/M_{\mathrm{Pl}}}$ simuliert. Das Skalarfeld-Minimum $\varphi_{\min}(\rho)$ zeigt einen klaren Phasenuebergang: fuer $\rho \lesssim 10^{-8}$ (kosmologisch) ist das Skalaron aktiv ($\varphi_{\min} \approx -5$), waehrend es fuer $\rho \gtrsim 10^{3}$ (Sonnensystem) vollstaendig gescreent ist ($\varphi_{\min} \to 0$). Die $\Gamma$-Konvergenz fuer $\varepsilon \to 0$ erzeugt eine scharfe Phasengrenze und bestaetigt, dass Screening als thermodynamischer Phasenuebergang entsteht, nicht als dynamischer Prozess. Das Screening-Verhaeltnis $\rho_{\mathrm{solar}}/\rho_{\mathrm{crit}} \approx 10^{11}$ gewaehrleistet eine starke Unterdrueckung der fuenften Kraft. Skript: \texttt{compute\_screening\_phase.py}.
\end{remark}

\begin{proposition}[Screening-Viabilit\"atskriterium]\label{prop:screening-viability}
Sei $\Phi_\varepsilon[\rho, \varphi]$ das gemeinsame Funktional aus Schritt~1 mit effektivem Potential $V_{\mathrm{eff}}(\varphi; \rho) = V(\varphi) + \rho\, A(\varphi)$. Definiere die Skalaronmasse
\begin{equation}
  m_s^2(\rho) := \frac{\partial^2 V_{\mathrm{eff}}}{\partial \varphi^2}\bigg|_{\varphi = \varphi_{\min}(\rho)}.
\end{equation}
Falls $A(\varphi) = e^{\varphi/M_{\mathrm{Pl}}}$ (konforme Kopplung), dann gilt:
\begin{enumerate}
\item[(i)] $m_s^2(\rho) \to m_0^2 > 0$ f\"ur $\rho \to 0$ (kosmologisches Regime: endliche Masse, Dunkle Energie aktiv);
\item[(ii)] $m_s^2(\rho) \sim \rho / M_{\mathrm{Pl}}^2$ f\"ur $\rho \to \infty$ (dichtes Regime: Masse w\"achst mit Dichte);
\item[(iii)] Die Yukawa-Unterdr\"uckungsl\"ange $\lambda_Y = 1/m_s$ erf\"ullt $\lambda_Y(\rho_\odot) \ll 1\,\mathrm{AU}$ f\"ur Sonnensystem-Dichte $\rho_\odot$, sodass $|\gamma_{\mathrm{PPN}} - 1| \lesssim (m_s \cdot r)^{-2} e^{-m_s r}$ mit der Cassini-Schranke $|\gamma_{\mathrm{PPN}} - 1| < 2{,}3 \times 10^{-5}$ vertr\"aglich ist.
\end{enumerate}
\end{proposition}

\begin{proof}[Beweisskizze]
(i) Bei $\rho = 0$: $V_{\mathrm{eff}} = V(\varphi)$, also $m_s^2 = V''(\varphi_{\min})$. F\"ur das Hu--Sawicki-Potential gilt $V''(\varphi_{\min}) = \Lambda_0 \cdot (2/3)/M_{\mathrm{Pl}}^2 > 0$.
(ii) Bei gro\ss em $\rho$: $\varphi_{\min} \to 0$ und $V_{\mathrm{eff}}'' \approx V''(0) + \rho A''(0) = V''(0) + \rho/M_{\mathrm{Pl}}^2$. Der $\rho$-abh\"angige Term dominiert.
(iii) F\"ur $\rho_\odot \sim 10^3$ (in nat\"urlichen Einheiten) gilt $m_s \sim \sqrt{\rho_\odot}/M_{\mathrm{Pl}} \gg H_0$, also $\lambda_Y \ll H_0^{-1}$. Bei $r = 1\,\mathrm{AU}$ ist der Yukawa-Faktor $e^{-m_s r}$ exponentiell klein.
\end{proof}

\begin{remark}[Screening als Phasen\"ubergang]\label{rem:screening-phase-transition}
Proposition~\ref{prop:screening-viability} zeigt, dass die $\Gamma$-konvergente Phasenseparation aus den Schritten~1--4 automatisch viables Screening produziert: die dichte Phase (Sonnensystem) hat $m_s \gg H_0$, was die f\"unfte Kraft unterdr\"uckt, w\"ahrend die verd\"unnte Phase (kosmologisch) $m_s \sim H_0$ hat, sodass Dunkle Energie die Beschleunigung antreiben kann. Anders als bei traditionellen Chameleon-Mechanismen ist dieses Screening eine \emph{Konsequenz} der Variationsstruktur, kein zus\"atzliches Postulat. Numerische Best\"atigung: \texttt{compute\_screening\_phase.py} (Bemerkung~\ref{rem:screening-numerics}).
\end{remark}

\subsubsection*{Schwellenaxiom: RG-Matching-Konsistenz}

\begin{axiom}[RG-PPN-Konsistenz -- RG-DE]\label{ax:rg-ppn}
Es existiert eine effektive Feldtheorie $S_{\mathrm{eff}}(\mu)$ mit RG-Skala $\mu$, sodass:
\begin{enumerate}
\item[(RG1)] \textbf{Eindeutiges Matching:} Die Parameter von $S_{\mathrm{eff}}$ werden bei einer Referenzskala $\mu_0$ durch physikalische Observablen ($H_0$, $\Omega_m$, $G_N$) fixiert.
\item[(RG2)] \textbf{Skalenstabilit\"at:} Vorhersagen f\"ur PPN-Parameter, insbesondere $\gamma_{\mathrm{PPN}}$, sind unabh\"angig von $\mu$ bis auf kontrollierte Korrekturen h\"oherer Ordnung: $|\gamma_{\mathrm{PPN}}(\mu) - \gamma_{\mathrm{PPN}}(\mu_0)| \leq \delta(\mu)$ mit $\delta \to 0$.
\item[(RG3)] \textbf{Limeskonsistenz:} Der Schwachfeld- (Newtonsche), Starkfeld- (relativistische) und kosmologische Limes kommutieren mit dem RG-Fluss.
\end{enumerate}
\end{axiom}

\begin{remark}[Der Matching-Leck-Scanner]\label{rem:matching-leak}
Um zu lokalisieren, wo die RG-PPN-Konsistenz im aktuellen Modell bricht:
\begin{enumerate}
\item \textbf{Parameterdrift:} Das lineare $R^2$-Modell liefert $\gamma_{\mathrm{PPN}} = 1/2$ auf Sonnensystem-Skalen (Bemerkung~\ref{sec:screening}), was Cassini verletzt ($|\gamma_{\mathrm{PPN}} - 1| < 2{,}3 \times 10^{-5}$). Die Skalaronmasse $m_s \ll H_0$ ist die Ursache.
\item \textbf{Feldredefinitions-Ambiguit\"at:} Die Jordan-Rahmen-$f(R)$- und Einstein-Rahmen-Skalar-Tensor-Formulierungen liefern identische Vorhersagen nur dann, wenn die konforme Transformation auf allen Skalen konsistent durchgef\"uhrt wird.
\item \textbf{Limes-Nichtkommutativit\"at:} Der \"Ubergang $\mu \to \mu_{\mathrm{solar}}$ und dann $R \to R_\odot$ kann sich von der umgekehrten Reihenfolge unterscheiden, falls der Screening-\"Ubergang nicht glatt ist.
\end{enumerate}
Die Hu--Sawicki-Parametrisierung ($f(R) = R - m^2 c_1 (R/m^2)^n / (c_2 (R/m^2)^n + 1)$) l\"ost Problem~(1), indem $m_s(\rho) \sim \sqrt{\rho}$ dichteabh\"angig gemacht wird (Proposition~\ref{prop:screening-viability}). Die quantitative Bedingung lautet: $n \geq 1$ und $|f_{R0}| \lesssim 10^{-6}$ zur Erf\"ullung der Cassini-Schranke.
\end{remark}

\begin{remark}[No-Go-Mechanismen]\label{rem:rg-nogo}
RG-PPN versagt, falls:
\begin{enumerate}
\item \textbf{Skalenabh\"angige Reparametrisierung:} Verschiedene Skalen erfordern verschiedene ,,nat\"urliche'' Parameter, was $\gamma_{\mathrm{PPN}}$ schemaabh\"angig macht -- ein EFT-Zusammenbruch.
\item \textbf{Doppelz\"ahlung:} Falls UV- und IR-Freiheitsgrade beide gez\"ahlt werden, wird $\Lambda_{\mathrm{eff}}$ \"ubergez\"ahlt und $\gamma_{\mathrm{PPN}}$ erbt die Inkonsistenz.
\item \textbf{Nicht-kanonische Limites:} Falls der Newtonsche oder post-Newtonsche Limes nicht eindeutig ist, ist jede PPN-Vorhersage physikalisch leer.
\end{enumerate}
Das $\Gamma$-konvergente Screening (Abschnitt~\ref{sec:screening-gamma}) adressiert Mechanismus~(1) durch eine variationelle (schemaunabh\"angige) Formulierung. Mechanismen~(2) und~(3) erfordern die explizite RG-Flussgleichung, die offen bleibt.
\end{remark}

\begin{remark}[Die EFT-Matching-Kette]\label{rem:eft-chain}
Die Konsistenz von Axiom~\ref{ax:rg-ppn} reduziert sich auf die Kommutativitaet eines einzigen Diagramms:
\[
  f(R)\text{-Parameter} \xrightarrow{\text{EFT-Matching}} m_s(\rho) \xrightarrow{\text{PPN-Limes}} \gamma_{\mathrm{PPN}},
\]
wobei jeder Pfeil unabhaengig von der Matching-Konvention sein muss (bis auf kontrollierte Korrekturen). Sonnensystem-Tests sind kein ,,externes Loch'', sondern der \emph{Konsistenz-Waechter}: Falls $\gamma_{\mathrm{PPN}}$ nur durch Reparametrisierung kompatibel gemacht werden kann, hat das Matching einen Defekt. Das aktuelle Modell hat $\gamma_{\mathrm{PPN}} = 1/2$ fuer lineares $R^2$; die Hu--Sawicki-Nichtlinearitaet ($n \geq 1$, $|f_{R0}| \leq 10^{-6}$) stellt $\gamma_{\mathrm{PPN}} \approx 1$ wieder her, indem $m_s$ dichteabhaengig wird (Proposition~\ref{prop:screening-viability}).
\end{remark}

\begin{lemma}[RG-PPN impliziert Beobachtungsviabilit\"at]\label{lem:rg-viability}
Falls Axiom~\ref{ax:rg-ppn} mit Hu--Sawicki-Parametern $n \geq 1$, $|f_{R0}| \leq 10^{-6}$ gilt, dann:
\begin{enumerate}
\item[(a)] $|\gamma_{\mathrm{PPN}} - 1| < 2{,}3 \times 10^{-5}$ (Cassini-kompatibel),
\item[(b)] $w_{\mathrm{DE}}(z=0) \in [-1{,}2;\, -0{,}8]$ (DESI DR2-kompatibel),
\item[(c)] BBN- und CMB-Bedingungen sind in f\"uhrender Ordnung erf\"ullt.
\end{enumerate}
\end{lemma}

\begin{remark}[Zwei-Skalen-$\gamma(R)$ f\"ur Chameleon-Screening]
\label{rem:two-scale-gamma}
Das Modell $f(R) = R + \gamma R^2$ erfordert $\gamma \sim 10^{60}$ in nat\"urlichen Einheiten, um $\Lambda_{\mathrm{eff}} \sim H_0^2$ zu reproduzieren. Dieser gro\ss e Wert ist im Sonnensystem problematisch, wo $R \gg H_0^2$ gilt. Eine laufende Kopplung $\gamma(R)$ l\"ost dies:
\[
\gamma(R) = \gamma_0 \cdot \left(\frac{R_0}{R}\right)^n, \quad n > 0,
\]
wobei $R_0 \sim H_0^2$ und $\gamma_0 \sim 10^{60}$. F\"ur $R \gg R_0$ (Sonnensystem) gilt $\gamma(R) \ll 1$ und die Skalaronmasse $m_s^2 \sim 1/(6\gamma(R))$ wird gro\ss, was die f\"unfte Kraft unterdr\"uckt (Chameleon-Mechanismus). F\"ur $R \sim R_0$ (kosmologisch) gilt $\gamma(R) \sim \gamma_0$ und das Modell reproduziert das beobachtete $\Lambda_{\mathrm{eff}}$. Das Hu--Sawicki-Modell~\cite{HuSawicki2007} liefert eine konkrete Realisierung mit $n = 1$.
\end{remark}

\begin{remark}[Viable $f(R)$-Klasse]
\label{rem:viable-fR}
Eine viable $f(R)$-Theorie muss drei Bedingungen gleichzeitig erf\"ullen:
\begin{enumerate}
\item \emph{Thermodynamisches Minimum:} Das Funktional $\Phi[\rho]$ besitzt ein stabiles de-Sitter-Minimum mit $\Lambda_{\mathrm{eff}} \sim H_0^2$ (gew\"ahrleistet durch (C1)--(C3) aus Theorem~\ref{thm:lambda-selection}).
\item \emph{PPN-Grenzen:} Der parametrisierte post-Newtonsche Parameter $\gamma_{\mathrm{PPN}} = 1 - 2m_s^{-2} R / (3 + 2m_s^{-2} R)$ erf\"ullt $|\gamma_{\mathrm{PPN}} - 1| < 2{,}3 \times 10^{-5}$ (Cassini-Schranke), was $m_s \gtrsim 10^{-3}$\,eV im Sonnensystem erfordert.
\item \emph{Keine Geister:} $f''(R) > 0$ (Skalaron ist nicht tachyonisch) und $f'(R) > 0$ (Graviton hat positive kinetische Energie) f\"ur alle $R > 0$.
\end{enumerate}
Die Hu--Sawicki-Klasse $f(R) = R - \mu^2 c_1 (R/\mu^2)^n / (1 + c_2 (R/\mu^2)^n)$ mit $n \geq 1$, $c_1/c_2 = 6\Omega_\Lambda/\Omega_m$ und $\mu^2 = H_0^2 \Omega_m$ erf\"ullt alle drei Bedingungen. Das FST-Funktional $\Phi[\rho]$, ausgewertet auf dieser Klasse, liefert $\Lambda_{\mathrm{eff}}$ konsistent mit Planck 2018 und DESI DR1.
\end{remark}

% ============================================================
\subsection{Das Resolventen-Vakuumproblem}\label{sec:resolvent-vacuum}

Die kosmologische Konstante ist fundamental ein
\emph{Resolventen-Problem zweiter Ordnung} im Sinne des universellen
Musters, das \"uber das gesamte FST-Programm hinweg identifiziert
wurde~\cite{FST-Unified2026,GeigerRH2026c}. Die dreistufige Hierarchie
manifestiert sich wie folgt:

\begin{enumerate}
  \item \textbf{F\"uhrender neutraler Beitrag ($O(1)$):}\; Die naive
        Vakuumenergiedichte $\rho_0 \sim M_{\mathrm{Pl}}^4$ ist der
        Beitrag f\"uhrender Ordnung. Er ist ,,neutral`` in dem Sinne,
        dass er den beobachteten Wert um $\sim 120$ Gr\"o\ss enordnungen
        massiv \"uberschreitet und keine n\"utzliche Einschr\"ankung der
        physikalischen kosmologischen Konstante liefert.
  \item \textbf{Entartung erster Ordnung ($O(1/L)$):}\;
        Quantenschleifenkorrekturen---Ein-Schleifen-, Zwei-Schleifen- und
        h\"ohere---erzeugen Beitr\"age zu $\rho_{\mathrm{vac}}$, die
        UV-divergent sind und Renormierung erfordern. Nach der
        Renormierung wird die Vakuumenergiedichte zum freien Parameter
        $\Lambda_{\mathrm{ren}}(\mu)$ auf jeder Energieskala $\mu$. Dies
        ist ein ,,entarteter`` Zustand: Der renormierte Wert kann beliebig
        sein, und kein Selektionsprinzip wirkt auf dieser Ordnung.
  \item \textbf{Resolvente zweiter Ordnung ($O(1/L^2)$):}\; Die
        Beitr\"age des topologischen Sektors---insbesondere das
        Zusammenspiel von de-Sitter-Horizontentropie, dem gravitativen
        R\"uckwirkungsterm~\eqref{eq:Vgrav} und der
        P\"oschl--Teller-S\"attigung~\eqref{eq:VPT}---erzeugen eine
        selbstkonsistente Selektion von $\rho_{\mathrm{CC}} \sim
        H_0^2\,M_{\mathrm{Pl}}^2 \sim 10^{-120}\,M_{\mathrm{Pl}}^4$.
        Das Freie-Energie-Funktional $\Phi(\rho)$
        (Theorem~\ref{thm:main}) kodiert diese
        Resolventenstruktur zweiter Ordnung: Seine strikte Konvexit\"at
        gew\"ahrleistet einen eindeutigen Minimierer, und die Skalierung
        des Minimierers wird durch das Gleichgewicht zwischen den
        entropischen und gravitativen Termen auf der IR-Skala
        $\mu = H_0$ bestimmt.
\end{enumerate}

\begin{remark}[Universelles Resolventen-Muster]\label{rem:universal-resolvent}
Das Resolventen-Vakuumproblem ist eine Instanz des
\emph{Resolventen-Dominanz-Musters zweiter Ordnung}, das \"uber alle f\"unf
Probleme im FST-Programm identifiziert wurde~\cite{FST-Unified2026}.
Im Fall der Dunklen Energie nimmt die dreistufige Hierarchie
folgende konkrete Gestalt an:

\begin{center}
\begin{tabular}{p{3.5cm}p{8.5cm}}
\hline
\textbf{Stufe} & \textbf{Dunkle-Energie-Instanziierung} \\
\hline
F\"uhrender neutraler Beitrag & Vakuumenergiedichte $\rho_0$ -- identisch f\"ur alle Feldkonfigurationen \\
1.\ Ordnung entartet & Schleifenkorrekturen entartet -- perturbative Strahlungskorrekturen k\"onnen kein eindeutiges $\Lambda$ selektieren \\
2.\ Ordnung entscheidet & $\Lambda_{\mathrm{eff}}$-Selektion durch entropisch-gravitatives Gleichgewicht bei $\mu = H_0$ \\
\hline
\end{tabular}
\end{center}

\noindent
Dieselbe dreistufige Struktur tritt in vier Begleitproblemen auf
(RH, Yang--Mills-Massel\"ucke, Navier--Stokes-Regularit\"at, turbulente Kaskade);
der detaillierte Vergleich findet sich in~\cite{FST-Unified2026}.
\end{remark}

\noindent
Diese Resolventen-Interpretation stellt das Problem der kosmologischen
Konstante auf dasselbe strukturelle Fundament wie die Riemannsche
Hypothese und die Yang--Mills-Massel\"ucke: In jedem Fall ist die
,,L\"osung`` auf erster Ordnung unsichtbar und tritt erst durch eine
Kopplung zweiter Ordnung zwischen entarteten und komplement\"aren
Spektralsektoren hervor. Der thermodynamische Selektionsmechanismus
des Theorems~\ref{thm:main} ist die Dunkle-Energie-Instanziierung
dieses universellen Musters.

\begin{remark}[Wasserstein-Konvexitaet und Talagrand-Schranke]
\label{rem:talagrand}
Das Funktional $\Phi[\rho]$ ist gleichmaessig konvex in der Wasserstein-2-Metrik auf dem Raum der Dichteprofile, mit Konvexitaetskonstante $\lambda_\Phi$, bestimmt durch den gravitativen Selbstwechselwirkungsterm. Nach dem Satz von Otto--Villani ist das de-Sitter-Vakuum $\rho^*$ das eindeutige Wasserstein-stabile Gleichgewicht: $W_2(\rho, \rho^*) \leq \sqrt{2 \Phi[\rho] / \lambda_\Phi}$. Dies liefert eine thermodynamische Stabilitaetsgarantie fuer die kosmologische Konstante, die unabhaengig vom UV-Cutoff ist.
\end{remark}

\subsection*{Beweisstatus-\"Ubersicht}

\begin{center}
\begin{tabular}{|l|l|}
\hline
\textbf{Komponente} & \textbf{Status} \\
\hline
\hline
\multicolumn{2}{|l|}{\textit{Bewiesen / hergeleitet:}} \\
\hline
$\Lambda_{\mathrm{eff}}$-Selektion aus freier Energie & Theorem \\
$\Lambda \sim H^2 / \log(M_{\mathrm{Pl}}/H)$-Skalierung & Hergeleitet \\
Saettigungsdynamik (tanh-Profil) & Proposition \\
$\Gamma$-Konvergenz des Screenings (Schritte 1--4) & Bewiesen \\
Screening-Viabilitaet ($m_s^2 \sim \rho$) & Proposition \\
\hline
\multicolumn{2}{|l|}{\textit{Numerisch bestaetigt:}} \\
\hline
DESI DR2 Kompatibilitaet ($0{,}39\sigma$) & Best-Fit $\kappa = 2{,}63$ \\
Phasenuebergang bei $\rho_{\mathrm{crit}} \sim 10^{-8}$ & Simulation \\
Screening-Verhaeltnis $\rho_\odot / \rho_c \sim 10^{11}$ & Simulation \\
\hline
\multicolumn{2}{|l|}{\textit{Offen / skizziert:}} \\
\hline
RG-Flussgleichung (UV-Vervollstaendigung) & Skizziert \\
Expliziter Hu--Sawicki-Parameterfit & Offen \\
CMB + BBN gemeinsame Einschraenkungen & Offen \\
\hline
\end{tabular}
\end{center}

% ============================================================
\section{Diskussion und Ausblick}\label{sec:discussion}

\subsection{Beziehung zum de-Sitter-Gleichgewicht}

Die de-Sitter-Temperatur $T_{\mathrm{dS}} = H/(2\pi)$ spielt eine
zentrale Rolle in unserem Funktional. Ein reines de-Sitter-Universum
mit konstantem~$H$ befindet sich im thermodynamischen Gleichgewicht
mit seinem kosmologischen Horizont, und unser Ergebnis kann als die
Aussage interpretiert werden, dass das beobachtete Universum sich
nahe diesem Gleichgewicht befindet. Abweichungen entstehen durch den
Materie- und Strahlungsgehalt, der das Minimum des Funktionals st\"ort.

\subsection{Einschr\"ankungen}

Mehrere Aspekte der Herleitung bed\"urfen weiterer Entwicklung:
\begin{enumerate}
  \item \textbf{Gravitative R\"uckwirkung:} Die funktionale Form
        von~$V_{\mathrm{grav}}$ wird nun durch zwei unabh\"angige
        Argumente gest\"utzt (Lemma~\ref{lem:minimality} und
        Satz~\ref{prop:Vgrav-from-fR}), doch die Herleitung aus
        der $f(R)$-Spurgleichung beruht auf der quasistatischen
        Sub-Horizont-N\"aherung. Eine vollst\"andig kovariante
        Herleitung aus der semiklassischen effektiven Wirkung
        (z.\,B.\ dem Pfadintegral des konformen Faktors) w\"urde die
        Grundlage verst\"arken.
  \item \textbf{UV-Renormierung und der Faktor-3-Diskrepanz:}
        Die Modensumme \"uber alle Standardmodell-Spezies und die
        ordnungsgem\"a\ss e Renormierung der UV-Divergenzen werden nur
        schematisch behandelt. Die aktuelle $\sim\!3\times$-Diskrepanz
        zwischen unserer Absch\"atzung
        ($\Leff \approx 3.8 \times 10^{-53}\;\mathrm{m}^{-2}$) und
        dem Planck-2018-Wert ($1.11 \times 10^{-52}\;\mathrm{m}^{-2}$)
        l\"asst sich auf drei Quellen zur\"uckf\"uhren:
        (i)~den $O(1)$-Koeffizienten~$c_0$ im R\"uckwirkungsterm,
        der durch die $f(R)$-Herleitung auf sub-Compton-Skalen um~$4/3$
        verst\"arkt wird; (ii)~den speziesabh\"angigen
        Multiplizit\"atsfaktor ($N_{\mathrm{eff}} \approx 28.75$
        Helizit\"atszust\"ande oberhalb $1\;\mathrm{MeV}$); und
        (iii)~das pr\"azise Renormierungsschema.  Eine vollst\"andige
        Berechnung, die alle drei Effekte einbezieht, w\"urde den
        Vorfaktor auf $O(1)$-Genauigkeit fixieren.
  \item \textbf{Vollst\"andige Stabilit\"atsanalyse:} Die
        de-Sitter-Stabilit\"at
        (Abschnitt~\ref{sec:self-consistent}) ist nur auf der
        Skalarfeld-Ebene etabliert. Eine vollst\"andige Analyse
        einschlie\ss lich Metrikst\"orungen im
        Mukhanov--Sasaki-Formalismus wird auf zuk\"unftige Arbeiten
        verschoben.
\end{enumerate}

\subsection{Renormierung und Status des Skalierungsgesetzes}
\label{sec:renorm}

Das Freie-Energie-Funktional~$\Phi(\rho)$ wie in~\eqref{eq:Phi}
definiert verwendet harte Impuls-Cutoffs $\LUV = M_{\mathrm{Pl}}$
und $\LIR = H_0$. Unter Standard-Renormierung
(Zetafunktion-, dimensionale oder adiabatische Subtraktion) erzeugt
die Nullpunktsenergie $\tfrac{1}{2}\omega(k)$ Divergenzen proportional
zu den lokalen geometrischen Invarianten $\sqrt{-g}$,
$\sqrt{-g}\,R$ und $\sqrt{-g}\,R^2$, die in Gegenterme f\"ur die
kosmologische Konstante, die Newtonsche Konstante bzw.\ h\"ohere
Kr\"ummungskopplungen absorbiert werden (vgl.~\cite{BirrellDavies1982}).
Nach der Renormierung ist die \emph{absolute} Vakuumenergiedichte keine
berechenbare Ausgabe, sondern ein renormierter Parameter
$\Lambda_{\mathrm{ren}}(\mu)$, der durch eine Anpassungsbedingung
fixiert wird.

Das Skalierungsgesetz $\Leff \sim H_0^2/\ln(M_{\mathrm{Pl}}/H_0)$
des Theorems~\ref{thm:main} sollte daher nicht als rohe
Modensummen-Vorhersage verstanden werden, sondern als eine
\emph{Renormierungsgruppen-Anpassungsaussage}: Das Funktional $\Phi$
wird auf der Infrarotskala $\mu = H$ ausgewertet, und der Logarithmus
$\ln(M_{\mathrm{Pl}}/H)$ entsteht aus dem Laufen von
$\Lambda_{\mathrm{ren}}(\mu)$ zwischen dem UV-Anpassungspunkt
$\mu = M_{\mathrm{Pl}}$ und der physikalischen Skala $\mu = H_0$.
In der Notation der Renormierungsgruppe:
\begin{equation}\label{eq:RG-running}
  \Lambda_{\mathrm{ren}}(H)
  \;=\; \Lambda_{\mathrm{ren}}(M_{\mathrm{Pl}})
  \;+\; \beta_\Lambda\,\ln\!\left(\frac{H}{M_{\mathrm{Pl}}}\right)
  \;+\; \cdots
\end{equation}
Die Betafunktion $\beta_\Lambda$ erh\"alt Beitr\"age sowohl von der
gravitativen R\"uckwirkung (kodiert in $V_{\mathrm{grav}}$) als auch
von den nichtlokalen Infraroteffekten des de-Sitter-Horizonts. Auf
einem quasi-de-Sitter-Hintergrund werten sich nichtlokale Terme der
schematischen Form $R\,\ln(\Box/\mu^2)\,R$ in der effektiven Wirkung
zu Logarithmen von $H/\mu$ aus, was einen physikalischen
(nicht-cutoff-basierten) Ursprung f\"ur den logarithmischen Faktor
liefert. In dieser Lesart ist der
R\"uckwirkungsterm~\eqref{eq:Vgrav} eine effektive, integrierte
Darstellung dieser nichtlokalen Beitr\"age, und die Skalierung
$\rho_\ast \sim H^2 M_{\mathrm{Pl}}^2/\ln(M_{\mathrm{Pl}}/H)$ ist
eine Aussage \"uber das infrarote Laufen und kein Cutoff-Artefakt.

Wir k\"onnen $\beta_\Lambda$ direkt aus dem Freie-Energie-Funktional
extrahieren, ohne auf QFT-Schleifenrechnungen zur\"uckgreifen zu m\"ussen.

\begin{proposition}[Thermodynamische $\beta$-Funktion]\label{prop:beta-Lambda}
Definiere die thermodynamische Antwortfunktion
\begin{equation}\label{eq:beta-thermo}
  \beta_\Lambda^{\mathrm{thermo}}
  \;:=\;
  \frac{\dd}{\dd\ln a}\,
  \frac{\partial\Phi}{\partial\rho}\bigg|_{\rho=\rho_\ast}.
\end{equation}
Auswertung am Minimum $\rho_\ast$ von~$\Phi$ ergibt
\begin{equation}\label{eq:beta-explicit}
  \beta_\Lambda^{\mathrm{thermo}}
  \;=\;
  -\,\frac{H_0^2\,M_{\mathrm{Pl}}^2}{\bigl[\ln(M_{\mathrm{Pl}}/H)\bigr]^2}
  \;\cdot\; \frac{\dd\ln H}{\dd\ln a},
\end{equation}
Auswertung in der Materie\"ara, wo $\dd\ln H/\dd\ln a = -3/2$:
\begin{equation}
  \beta_\Lambda^{\mathrm{thermo}}\Big|_{\mathrm{Materie}}
  \;=\;
  +\,\frac{3}{2}\;\frac{H_0^2\,M_{\mathrm{Pl}}^2}
  {\bigl[\ln(M_{\mathrm{Pl}}/H)\bigr]^2}
  \;>\; 0.
\end{equation}
Das positive Vorzeichen bedeutet, dass $\Leff$ zunimmt, w\"ahrend das
Universum expandiert und $H$ abnimmt---konsistent mit dem \"Ubergang
von Verz\"ogerung zu Beschleunigung.
\end{proposition}

\begin{proof}
Die renormierte Stationarit\"atsbedingung (siehe
Abschnitt~\ref{sec:renorm}) definiert implizit
$\rho_\ast = \rho_\ast(H)$, da $\Phi$ \"uber die
de-Sitter-Temperatur $T_{\mathrm{dS}} = H/(2\pi)$ und den
IR-Cutoff $\LIR = H$ von~$H$ abh\"angt. Nach dem Satz \"uber
implizite Funktionen, Differentiation von
$\dd\Phi/\dd\rho\big|_{\rho_\ast} = 0$ nach $\ln a$ bei festem
UV-Cutoff:
\[
  \frac{\partial^2\Phi}{\partial\rho^2}\bigg|_{\rho_\ast}
  \cdot \frac{\dd\rho_\ast}{\dd\ln a}
  \;=\;
  -\frac{\partial^2\Phi}{\partial\rho\,\partial H}
  \bigg|_{\rho_\ast}
  \cdot \frac{\dd H}{\dd\ln a}.
\]
Die dominante $H$-Abh\"angigkeit geht \"uber
$B \propto T_{\mathrm{dS}} \propto H$ ein. In der renormierten
Skalierung
$\rho_\ast^{\mathrm{ren}} \sim H^2 M_{\mathrm{Pl}}^2/\ln(M_{\mathrm{Pl}}/H)$
ergibt die logarithmische Ableitung
$\dd\ln\rho_\ast/\dd\ln H \approx 2 + 1/\ln(M_{\mathrm{Pl}}/H)
\approx 2$ den Zusammenhang
$\dd\rho_\ast/\dd\ln a \approx 2\rho_\ast\,\dd\ln H/\dd\ln a$.
Die thermodynamische Betafunktion~\eqref{eq:beta-thermo} wertet sich
damit zu~\eqref{eq:beta-explicit} aus.
In der Materie\"ara gilt $H \propto a^{-3/2}$, also
$\dd\ln H/\dd\ln a = -3/2$, und die beiden Minuszeichen (eines aus
der Formel, eines aus $\dd\ln H/\dd\ln a < 0$) ergeben zusammen
$\beta_\Lambda^{\mathrm{thermo}} > 0$.
\end{proof}

\begin{remark}
Satz~\ref{prop:beta-Lambda} schlie\ss t die RG-L\"ucke (L3), die in
der Reviewanalyse identifiziert wurde: $\beta_\Lambda$ ist nun eine
explizite, berechenbare Funktion des Hubble-Parameters, vollst\"andig
aus dem thermodynamischen Funktional~$\Phi$ abgeleitet, ohne
QFT-Schleifenrechnungen. Vorzeichen und Gr\"o\ss enordnung sind
konsistent mit dem
Skalar-Tensor-Lagrangian~\eqref{eq:Lagrangian}: Das
Skalaron-vermittelte Laufen
reproduziert~\eqref{eq:beta-explicit} auf Tree-Level.
\end{remark}

\begin{remark}[Semiklassische Effektivwirkung]
\label{rem:semiclassical}
Die Ein-Schleifen-Effektivwirkung f\"ur $f(R) = R + \gamma R^2$-Gravitation hat die Form
\[
\Gamma_{\mathrm{1-loop}} = \int d^4x \sqrt{-g} \left[ \frac{R}{16\pi G} + \gamma R^2 + \frac{1}{(4\pi)^2} \left( c_1 R^2 \log\frac{R}{\mu^2} + c_2 R_{\mu\nu} R^{\mu\nu} \log\frac{R}{\mu^2} \right) \right],
\]
wobei $c_1, c_2$ schema-abh\"angige Konstanten und $\mu$ die Renormierungsskala ist. Die nichtlokalen Terme $R \log(\Box/\mu^2) R$ entstehen aus der W\"arme\-kern-Entwicklung und kodieren das Laufen von $\gamma$ mit der Skala. Die Wahl $\mu^2 = H_0^2$ liefert den kosmologischen Wert $\gamma_0 \sim 10^{60}$; die Wahl $\mu^2 = R_{\odot}$ (solare Kr\"ummung) ergibt $\gamma(R_{\odot}) \ll 1$, was eine feldtheoretische Begr\"undung f\"ur den Zwei-Skalen-Ansatz $\gamma(R)$ liefert.
\end{remark}

\begin{remark}[Betafunktion f\"ur $\Lambda_{\mathrm{eff}}$]
\label{rem:beta-lambda}
Im Wilsonschen Rahmen der effektiven Feldtheorie l\"auft die kosmologische Konstante mit der RG-Skala $k$:
\[
\beta_\Lambda = k \frac{d\Lambda(k)}{dk} = \frac{k^4}{16\pi^2} \left( \sum_i (-1)^{F_i} n_i - \frac{\Lambda(k)}{k^2} \right),
\]
wobei die Summe \"uber Teilchenspezies mit Spin-Statistik-Faktor $(-1)^{F_i}$ und Multiplizit\"at $n_i$ l\"auft. Das FST-Funktional $\Phi[\rho]$ liefert die Anpassungsbedingung auf der IR-Skala $k = H_0$: $\Lambda(H_0) = \Lambda_{\mathrm{eff}} \sim H_0^2$. Die UV-IR-Verbindung wird durch die Skalaronmasse $m_s$ vermittelt, die die \"Ubergangsskala zwischen dem laufenden Regime ($k > m_s$) und dem eingefrorenen Regime ($k < m_s$) festlegt.
\end{remark}

\subsection{Ein-Schleifen-Robustheit}\label{sec:one-loop}

Der Skalar-Tensor-Lagrangian~\eqref{eq:Lagrangian} enth\"alt die
$R^2$-Kopplung~$\gamma$, die einen zus\"atzlichen skalaren
Freiheitsgrad (das Skalaron) mit der Masse
$m_s^2 = 1/(6\gamma)$ einf\"uhrt. Sein
Ein-Schleifen-Beitrag zur Vakuumenergiedichte ist von der
Gr\"o\ss enordnung
\begin{equation}\label{eq:delta-rho-scalaron}
  \delta\rho_s
  \;\sim\; \frac{m_s^4}{64\pi^2}
  \;=\; \frac{1}{64\pi^2\,(6\gamma)^2}.
\end{equation}
Damit das Tree-Level-Minimum $\rho_\ast \sim 3M_{\mathrm{Pl}}^2 H_0^2$
nicht \"uberschattet wird, fordern wir $\delta\rho_s \ll \rho_\ast$,
was die parametrische Bedingung ergibt
\begin{equation}\label{eq:gamma-bound}
  \gamma \;\gg\; \frac{1}{M_{\mathrm{Pl}}\,H_0}
  \;\sim\; 10^{60}\;\text{in Planck-Einheiten}.
\end{equation}
F\"ur solch gro\ss e Werte von~$\gamma$ wird das Skalaron extrem leicht
($m_s \ll H_0$) und entkoppelt effektiv von der Sp\"atzeitdynamik.
Alternativ kann man $\rho_\ast$ als renormierten Anpassungswert
interpretieren, wobei $\delta\rho_s$ in die
Renormierungsbedingung f\"ur~$\Lambda_{\mathrm{ren}}$ absorbiert wird;
dies ist die Standardsichtweise in quadratischer
Gravitation~\cite{SotiriouFaraoni2010}.

Standardmodellfelder koppeln an $\phi$ ausschlie\ss lich \"uber die
Gravitation (es gibt keine direkten Yukawa- oder Portal-Kopplungen
im Lagrangian~\eqref{eq:Lagrangian}). Folglich sind
materieinduzierte Korrekturen zur Skalarfeldmasse
Planck-unterdr\"uckt:
$\delta m_\phi^2 \sim m_{\mathrm{heavy}}^2/(16\pi^2 M_{\mathrm{Pl}}^2)$,
was die ultraleichte Quintessenzmasse $m_\phi \sim H_0$ ohne
Feinabstimmung bewahrt. Falls jedoch nichtgravitative Kopplungen
vorhanden w\"aren, w\"urden radiative Korrekturen generisch
$\delta m_\phi^2 \sim y^2 m_{\mathrm{heavy}}^2/(16\pi^2)$ erzeugen,
was die technische Nat\"urlichkeit des P\"oschl--Teller-Potentials
verletzen w\"urde. Die Abwesenheit solcher Kopplungen ist daher eine
strukturelle Anforderung des Modells, konsistent mit dem rein
geometrischen Ursprung des Skalarfeldes im $R^2$-Sektor.

\subsection{Renormierungsgruppen-Interpretation von $\Phi$}
\label{sec:rg-interpretation}

Das Freie-Energie-Funktional $\Phi[\rho]$ besitzt eine nat\"urliche
Interpretation als Callan--Symanzik-Fluss zwischen UV- und
IR-Fixpunkten.

\begin{definition}[RG-Fluss von $\Phi$]
\label{def:rg-flow-phi}
Definiere die laufende freie Energie bei Skala~$\mu$:
\[
  \Phi(\rho;\,\mu)
  \;=\;
  \int d^{3}x\,\biggl[
    \rho\,\log\frac{\rho}{\rho_{0}(\mu)}
    - \rho + \rho_{0}(\mu)
    + \frac{G(\mu)}{2}\,\frac{\rho^{2}}{k_{\mathrm{IR}}^{2}(\mu)}
  \biggr],
\]
wobei $\rho_{0}(\mu)$, $G(\mu)$ und $k_{\mathrm{IR}}(\mu)$ mit der
RG-Skala~$\mu$ laufen.  Die Callan--Symanzik-Gleichung lautet
\[
  \biggl(
    \mu\,\frac{\partial}{\partial\mu}
    + \beta_{G}\,\frac{\partial}{\partial G}
    + \beta_{k}\,\frac{\partial}{\partial k_{\mathrm{IR}}}
    + \gamma_{\rho}\,\rho\,\frac{\partial}{\partial\rho}
  \biggr)\Phi = 0.
\]
\end{definition}

\begin{proposition}[Fixpunktstruktur]
\label{prop:rg-fixed-points}
Der RG-Fluss besitzt zwei Fixpunkte:
\begin{enumerate}
\item \emph{UV-Fixpunkt} ($\mu \to M_{\mathrm{Pl}}$):
  $\rho_{0} \sim M_{\mathrm{Pl}}^{4}/(8\pi G)$,
  $k_{\mathrm{IR}} \sim M_{\mathrm{Pl}}$.
  Die freie Energie wird vom Modenz\"ahlterm dominiert:
  $\Phi \sim M_{\mathrm{Pl}}^{4}\cdot\mathrm{vol}(M)$.
  Dies ist das Problem der kosmologischen Konstante in seiner
  urspr\"unglichen Form.
\item \emph{IR-Fixpunkt} ($\mu \to H_{0}$):
  $\rho^{*} = \rho_{0}(H_{0}) \sim H_{0}^{2}/(8\pi G)$,
  $k_{\mathrm{IR}} \sim H_{0}$.
  Die freie Energie wird am de-Sitter-Vakuum minimiert:
  $\Phi[\rho^{*}] = 0$.
  Die effektive kosmologische Konstante
  $\Lambda_{\mathrm{eff}} = 8\pi G\,\rho^{*}$ wird durch die
  IR-Skala bestimmt.
\end{enumerate}
Der Fluss von UV nach IR ist monoton:
$\mu\,\frac{d}{d\mu}\,\Phi[\rho^{*}(\mu)] \leq 0$
(die freie Energie nimmt unter Vergr\"oberung ab).
Dies ist das Analogon von Zamolodchikovs $c$-Theorem f\"ur die
kosmologische Konstante.
\end{proposition}

\begin{proof}[Beweisskizze]
Die Monotonie folgt aus der Konvexit\"at der KL-Divergenz: Bei jedem
RG-Schritt erf\"ullt die vergr\"oberte Dichte
$\rho_{k+1} = \mathcal{R}_{k}\,\rho_{k}$ die Ungleichung
$D_{\mathrm{KL}}(\rho_{k+1}\|\rho_{0})
 \leq D_{\mathrm{KL}}(\rho_{k}\|\rho_{0})$
gem\"a\ss{} der Datenverarbeitungsungleichung.  Der Gravitationsterm
$G\rho^{2}/k_{\mathrm{IR}}^{2}$ ist eine relevante St\"orung im
RG-Sinne (Dimension~2 in 4D), die zum IR hin w\"achst und den
Fixpunkt stabilisiert.
\end{proof}

\begin{remark}[Verbindung zur Resolventenarchitektur]
\label{rem:rg-resolvent}
Der RG-Fluss von $\Phi$ weist die Pattern-A-Resolventenstruktur auf:
\begin{itemize}
\item \emph{F\"uhrender neutraler Term:}
  Am UV-Fixpunkt ist $\Phi \sim M_{\mathrm{Pl}}^{4}$
  skalenunabh\"angig (nullte Ordnung).
\item \emph{Entartung erster Ordnung:}
  Die Ein-Schleifen-Korrektur
  $\delta\Phi \sim \log(\mu/M_{\mathrm{Pl}})$ ist eine logarithmische
  Korrektur, die den UV-Fixpunkt nicht bricht.
\item \emph{Dominanz zweiter Ordnung:}
  Die gravitative Selbstwechselwirkung
  $G\rho^{2}/k_{\mathrm{IR}}^{2}$ ist die relevante St\"orung, die
  den Fluss zum IR-Fixpunkt treibt und $\Lambda_{\mathrm{eff}}$
  bestimmt.
\end{itemize}
Dies ist strukturell identisch mit der Resolventenhierarchie in den
Begleitpapieren zu Yang--Mills (Poincar\'e-Konstante unter RG) und
Navier--Stokes (Gronwall-Konstante unter Galerkin-Trunkierung).
\end{remark}

\subsection{Vorhersagen: Langsame Evolution von $\Leff$}

Falls die effektive kosmologische Konstante wie
$\Leff \sim H^{2}/\ln(M_{\mathrm{Pl}}/H)$ evolviert, dann ist $\Leff$
\emph{nicht} exakt konstant, sondern nimmt langsam ab, w\"ahrend das
Universum expandiert und~$H$ abnimmt. Der resultierende intrinsische
Zustandsgleichungsparameter der Dunklen Energie erf\"ullt
$w_{\mathrm{DE}} = -1 + \varepsilon$ mit
\begin{equation}
  \varepsilon \;\sim\;
  \frac{2}{\ln(M_{\mathrm{Pl}}/H_0)}
  \;\approx\; 0.014,
\end{equation}
was mit den aktuellen Beschr\"ankungen
($w = -1.03 \pm 0.03$, Planck~2018) und den
DESI-DR2-Ergebnissen~\cite{DESI2025} konsistent ist, die
$2$--$4\sigma$-Evidenz f\"ur dynamische Dunkle Energie ($w \neq -1$)
zeigen. Surveys der n\"achsten Generation (Euclid,
DESI-Gesamtsurvey, Vera Rubin Observatory) werden diese
$\varepsilon \approx 0.014$-Abweichung mit der erforderlichen
Pr\"azision untersuchen.

\subsection{Thermodynamische Selektion der Vakuumenergie}\label{sec:selection}

Das Variationsergebnis des Theorems~\ref{thm:main} l\"asst eine komplement\"are thermodynamische Interpretation \"uber das Prinzip der maximalen Entropieproduktion (MEPP)~\cite{MartyushevSeleznev2006} zu. Unter allen kinematisch zul\"assigen Vakuumenergiedichten ist der realisierte Wert~$\rho_\ast$ derjenige, der die Rate der Entropieproduktion~$\dot{S}$ \"uber kosmologische Zeitskalen maximiert. Dieses Selektionskriterium ist konsistent mit unserer Variationsschranke: Das Minimum des Freie-Energie-Funktionals~$\Phi(\rho)$ f\"allt mit dem Maximum von~$\dot{S}$ zusammen, weil der Entropieterm in~$\Phi$ als thermodynamische Strafe wirkt, die sowohl \"uberm\"a\ss ig gro\ss e als auch verschwindend kleine Vakuumenergien unterdr\"uckt. Am Sattelpunkt strahlt der de-Sitter-Horizont bei der Temperatur $T_{\mathrm{dS}} = H/(2\pi)$, und die zugeh\"orige Gibbons--Hawking-Entropie $S_{\mathrm{dS}} = \pi c^5/(G\hbar\,H^2)$ erreicht ein quasistation\"ares Maximum---das Universum expandiert gerade schnell genug, um das zug\"angliche Phasenraumvolumen zu maximieren.

Aus dieser Perspektive ist die beschleunigte Expansion kein Zufall, sondern ein thermodynamischer Imperativ. Sobald die stellare Nukleosynthese die verf\"ugbare baryonische freie Energie ersch\"opft hat, wird der dominante Kanal f\"ur Entropieproduktion die kosmologische Expansion selbst. Die Vakuumenergiedichte~$\rho_\ast \sim H_0^2/G$ ist pr\"azise der Wert, bei dem die Entropieproduktionsrate des Horizonts maximiert wird, unter der Nebenbedingung, dass~$\rho$ selbstkonsistent durch den gravitativen R\"uckwirkungsterm~\eqref{eq:Vgrav} aufrechterhalten werden muss. Dieser thermodynamische Selektionsmechanismus bietet einen physikalischen Grund, \emph{warum} die kosmologische Konstante den Wert annimmt, den sie hat, und erg\"anzt das variationelle \emph{Wie}, das durch das Freie-Energie-Funktional bereitgestellt wird.

In dieser Lesart l\"ost sich das Koinzidenzproblem vollst\"andig auf: Die Vakuumenergie folgt der Materiedichte nicht zuf\"allig, sondern weil beide von derselben entropiemaximierenden Dynamik gesteuert werden. Der Infrarot-Cutoff $\LIR = H$ koppelt den Vakuumsektor an die Expansionsgeschichte und stellt sicher, dass sich $\rho_\ast$ adiabatisch anpasst, w\"ahrend das Universum von der Materiedominanz zur Dunkle-Energie-Dominanz \"ubergeht.

\begin{remark}[2D-Parameterscan als quantitatives Koinzidenz-Analogon]
\label{rem:2d-scan-coincidence}
Der zweidimensionale Parameterscan \"uber $(\alpha, \gamma)$ aus der
begleitenden quantitativen Studie~\cite{FST-Unified2026} liefert eine
konkrete Realisierung der Koinzidenzproblem-Aufl\"osung.
Das Freie-Energie-Funktional $\Phi[\rho]$ weist in der
$(\alpha, \gamma)$-Ebene ein \emph{scharfes Minimum} auf, das
automatisch die beobachtete Gr\"o\ss enordnung der effektiven
kosmologischen Konstante~$\Lambda_{\mathrm{eff}}$ selektiert.
Diese Sch\"arfe wird nicht durch Fine-Tuning aufgepr\"agt, sondern
entsteht aus dem Zusammenspiel der Entropie-Strafe (logarithmisch,
bevorzugt kleines~$\rho$) und der gravitativen R\"uckwirkung
(quadratisch, bestraft verschwindendes~$\rho$). Der 2D-Scan erhebt
die Koinzidenzproblem-Aufl\"osung damit von einem qualitativen
Argument (``$\Lambda_{\mathrm{eff}}$ folgt~$H^2$'') zu einer
quantitativen Demonstration: Die Funktionallandschaft selbst zwingt
$\rho_\ast$ in ein enges Tal, dessen Tiefe durch $H_0^2/G$
bestimmt ist.
\end{remark}

\begin{theorem}[Thermodynamische Selektion von $\Lambda_{\mathrm{eff}}$]
\label{thm:lambda-selection}
Sei $\Phi[\rho]$ ein Freie-Energie-Funktional auf kosmologischen Dichteprofilen mit folgenden Eigenschaften:
\begin{enumerate}
\item[\textup{(C1)}] \emph{Allgemeine Kovarianz:} $\Phi$ ist invariant unter Koordinatentransformationen.
\item[\textup{(C2)}] \emph{Minkowski-Subtraktion:} Die nackte kosmologische Konstante $\Lambda_0 \sim M_{\mathrm{Pl}}^4$ wird durch Renormierung absorbiert: $\Phi[\rho] = \Phi_{\mathrm{ren}}[\rho - \rho_0]$, wobei $\rho_0$ das Minkowski-Vakuum ist.
\item[\textup{(C3)}] \emph{Stabilit\"at:} $\Phi$ besitzt einen eindeutigen Minimierer $\rho^*$ mit $\nabla^2 \Phi[\rho^*] > 0$.
\end{enumerate}
Dann erf\"ullt die effektive kosmologische Konstante $\Lambda_{\mathrm{eff}} = 8\pi G \rho^* \sim H_0^2$ und wird durch die Infrarotskala $H_0$ bestimmt, nicht durch die Ultraviolettskala $M_{\mathrm{Pl}}$.
\end{theorem}

\begin{proof}[Beweisskizze]
(C1) schr\"ankt $\Phi$ darauf ein, nur von skalaren Invarianten von $\rho$ abzuh\"angen.
(C2) entfernt den $M_{\mathrm{Pl}}^4$-Beitrag und l\"asst nur IR-Skalen-Terme \"ubrig.
(C3) selektiert das eindeutige Minimum, das durch Dimensionsanalyse in der
renormierten Theorie wie $H_0^2 / (8\pi G)$ skalieren muss.
Die zentrale Einsicht ist, dass (C2) kein Fine-Tuning darstellt, sondern eine
physikalische Renormierungsbedingung: Das Minkowski-Vakuum hat per Definition
eine kosmologische Konstante von null.
\end{proof}

\subsection{Falsifizierbare Vorhersagen}\label{sec:predictions-falsifiable}

Das Skalar-Tensor-Rahmenwerk aus Abschnitt~\ref{sec:scalar-tensor}
erzeugt eine Reihe quantitativer Vorhersagen, die es sowohl von
$\Lambda$CDM als auch von generischen modifizierten
Gravitationsmodellen unterscheiden. Wir listen die wichtigsten
Beobachtungssignaturen auf:

\begin{enumerate}
  \item \textbf{Gravitativer Slip-Parameter.}\;
        Der Skalar-Tensor-Lagrangian~\eqref{eq:Lagrangian} modifiziert die
        zwei Gravitationspotentiale $\Psi$ und $\Phi_{\mathrm{N}}$ in der
        konformen Newtonschen Eichung
        $\mathrm{d}s^{2} = -(1+2\Psi)\,\mathrm{d}t^{2}
        + a^{2}(1+2\Phi_{\mathrm{N}})\,\mathrm{d}\mathbf{x}^{2}$.
        In der quasistatischen Sub-Horizont-N\"aherung lauten die
        modifizierten Poissongleichungen f\"ur den
        $R + \gamma R^{2}$-Sektor~\cite{HuSawicki2007,SotiriouFaraoni2010}:
        \begin{align}
          k^{2}\,\Phi_{\mathrm{N}} &= -4\pi G\,a^{2}\,\bar{\rho}\,\delta
          \;\biggl[1 + \frac{1}{3}\,
          \frac{k^{2}/a^{2}}{k^{2}/a^{2} + m_{s}^{2}}\biggr],
          \label{eq:Poisson-Phi}\\[4pt]
          k^{2}\,\Psi &= -4\pi G\,a^{2}\,\bar{\rho}\,\delta
          \;\biggl[1 + \frac{2}{3}\,
          \frac{k^{2}/a^{2}}{k^{2}/a^{2} + m_{s}^{2}}\biggr],
          \label{eq:Poisson-Psi}
        \end{align}
        wobei $m_{s}^{2} = 1/(6\gamma)$ das Quadrat der Skalaronmasse ist.
        Der gravitative Slip-Parameter folgt unmittelbar:
        \begin{equation}\label{eq:grav-slip}
          \eta(k) \;\equiv\; \frac{\Psi}{\Phi_{\mathrm{N}}}
          \;=\;
          \frac{1 + \dfrac{2}{3}\,
          \dfrac{k^{2}/a^{2}}{k^{2}/a^{2} + m_{s}^{2}}}
          {1 + \dfrac{1}{3}\,
          \dfrac{k^{2}/a^{2}}{k^{2}/a^{2} + m_{s}^{2}}}
          \;\xrightarrow{\;k \gg a\,m_{s}\;}
          \;\frac{1 + \tfrac{2}{3}}{1 + \tfrac{1}{3}}
          \;=\; \frac{5/3}{4/3}
          \;=\; \frac{5}{4}.
        \end{equation}
        Im $\Lambda$CDM gilt $\eta = 1$ exakt. Der \"Ubergang findet bei
        $k_{\mathrm{cross}} \approx 0.86\,a\,m_{s}$ statt, wobei
        $\eta = 1.125$. F\"ur eine Skalaronmasse $m_{s} \sim H_{0}$
        (wie durch die P\"oschl--Teller-Dynamik nahegelegt) liegt die
        Compton-Wellenl\"ange $\lambda_{C} \sim c/H_{0}$ nahe dem
        Hubble-Radius, sodass die von Euclid untersuchten Skalen
        ($k \sim 0.01$--$0.2\;h\,\mathrm{Mpc}^{-1}$) tief im
        sub-Compton-Regime liegen, wo $\eta \approx 5/4$. Diese
        25\%-Abweichung von $\Lambda$CDM liegt innerhalb der
        Empfindlichkeit der Schwerkraftlinsen- und
        Galaxien-Clustering-Messungen von
        Euclid~\cite{HuSawicki2007}. Entscheidend ist, dass rein
        geometrische Alternativen wie die
        Finsler-Gravitation~\cite{FinslerJCAP2025}
        $\eta \neq 5/4$ vorhersagen, was den gravitativen Slip zu
        einer \emph{modelldiskriminierenden Observablen} zwischen dem
        CRM, $\Lambda$CDM und Finsler-Ans\"atzen macht.

  \item \textbf{Linsenparameter.}\;
        Der effektive Linsenparameter, definiert als
        $\Sigma(k) \equiv (\Phi_{\mathrm{N}} + \Psi)/
        (2\Phi_{\mathrm{N}})$, berechnet sich zu
        \begin{equation}\label{eq:Sigma}
          \Sigma(k) \;=\;
          \frac{2 + f(k)}{2 + \tfrac{2}{3}\,f(k)},
          \qquad
          f(k) \;:=\;
          \frac{k^{2}/a^{2}}{k^{2}/a^{2} + m_{s}^{2}}.
        \end{equation}
        Im GR-Limes ($f \to 0$, d.\,h.\ $k \ll a\,m_s$):
        $\Sigma \to 1$.  Im sub-Compton-Regime ($f \to 1$, d.\,h.\
        $k \gg a\,m_s$): $\Sigma \to 9/8 = 1{,}125$.  Die
        skalenabh\"angige Abweichung von $\Sigma = 1$ ist eine
        weitere unterscheidende Signatur gegen\"uber $\Lambda$CDM
        (wo $\Sigma = 1$ exakt gilt) und ist durch
        Gravitationslinsen-Surveys wie
        Euclid~\cite{SotiriouFaraoni2010} testbar.

  \item \textbf{Skalarfeldoszillationen in $H(a)$.}\;
        Zu sp\"aten Zeiten ($a \gtrsim 2$) oszilliert das
        Skalarfeld~$\phi$ um das Minimum von~$V_{\mathrm{PT}}$ mit
        der Frequenz $\omega_{\phi} = \sqrt{2V_{0}}/\phi_{0}$. Diese
        Oszillationen pr\"agen eine charakteristische Modulation auf den
        Hubble-Parameter:
        \begin{equation}
          \frac{\delta H}{H}
          \;\sim\; \frac{\Phi_{0}}{2}\;
          \mathrm{sech}^{2}\!\bigl[\kappa(a - a_{\mathrm{trans}})\bigr]\;
          \cos(\omega_{\phi}\,t),
        \end{equation}
        mit einer Amplitude der Gr\"o\ss enordnung
        $10^{-3}$--$10^{-2}$ bei $a \approx 2$. Zuk\"unftige
        Pr\"azisionsmessungen von $H(z)$ mittels
        Gravitationswellen-Standardsirenen k\"onnten diese
        oszillatorische Signatur detektieren.

  \item \textbf{Verbindung zur MOND-Beschleunigungsskala.}\;
        Die charakteristische Beschleunigungsskala des Skalarfeldes ist
        \begin{equation}\label{eq:a0-MOND}
          a_{0} \;=\; \frac{c\,H_{0}}{2\pi}
          \;\approx\; 1.0 \times 10^{-10}\;\mathrm{m\,s^{-2}},
        \end{equation}
        was mit der empirischen MOND-Beschleunigung
        $a_{0}^{\mathrm{MOND}} \approx 1.2 \times
        10^{-10}\;\mathrm{m\,s^{-2}}$ innerhalb von 20\%
        \"ubereinstimmt. Im vorliegenden Rahmenwerk ist dies kein
        Zufall: Dasselbe P\"oschl--Teller-Skalarfeld, das die
        Sp\"atzeitbeschleunigung antreibt, setzt auch die Skala, bei
        der die Gravitationsdynamik vom Newtonschen ins tiefe
        MOND-Regime \"ubergeht, \"uber die de-Sitter-Temperatur
        $T_{\mathrm{dS}} = \hbar\,a_{0}/(2\pi c\,k_{\mathrm{B}})$.

  \item \textbf{Resolventen-Vakuum-Vorhersage f\"ur $\eta$.}\;
        Das Resolventen-Rahmenwerk zweiter Ordnung
        (Abschnitt~\ref{sec:resolvent-vacuum}) liefert eine unabh\"angige
        Herleitung von $\eta = 5/4$: Das Gleichgewicht zwischen dem
        entropischen Term (de-Sitter-Temperatur) und der gravitativen
        R\"uckwirkung in $\Phi(\rho)$ fixiert das Verh\"altnis der
        modifizierten Poisson-Potentiale auf sub-Compton-Skalen und
        ergibt $\eta = 5/4$ als notwendige Konsequenz der
        Resolventenstruktur und nicht als Parameteranpassung.
        Entscheidend ist, dass
        Finsler-Gravitationsmodelle~\cite{FinslerJCAP2025}
        $\eta \neq 5/4$ vorhersagen, da ihre rein geometrische
        Modifikation der Friedmann-Gleichungen \"uber einen anderen
        Mechanismus wirkt (anisotrope Raumzeit-Metrik statt
        Skalar-Tensor-Kopplung). Dies macht $\eta$ zu einer scharfen
        \emph{modelldiskriminierenden Observablen}:
        \begin{itemize}
          \item \textbf{CRM/Resolvente:} $\eta = 5/4$ bei $k \gg a\,m_s$.
          \item \textbf{Finsler:} $\eta \neq 5/4$ (generisch $\eta > 1$,
                aber nicht $5/4$).
          \item \textbf{$\Lambda$CDM:} $\eta = 1$ exakt.
        \end{itemize}
        Der DESI-Gesamtsurvey (erwartet 2028) und Euclid~DR1 (erwartet
        2027) werden $\eta$ bei $z < 0.5$ mit
        $\sim 2\%$ Pr\"azision messen. Eine $2\sigma$-Abweichung von
        $5/4$ w\"urde das Resolventen-Rahmenwerk zweiter Ordnung f\"ur
        Dunkle Energie falsifizieren; Konvergenz zu $5/4$ w\"urde
        gleichzeitig sowohl $\Lambda$CDM als auch
        Finsler-Alternativen ausschlie\ss en.
\end{enumerate}

Jede dieser Vorhersagen ist mit Daten falsifizierbar, die innerhalb des
n\"achsten Jahrzehnts erwartet werden (Euclid DR1: 2027; LISA: 2030er;
DESI-Gesamtsurvey: 2028). Eine einzige best\"atigte Abweichung von
$\Lambda$CDM entlang eines dieser Kan\"ale w\"urde Evidenz f\"ur den
Skalar-Tensor-Mechanismus liefern; ein best\"atigtes $\eta = 1$ auf
sub-Compton-Skalen w\"urde ihn ausschlie\ss en.

\begin{remark}[Hu--Sawicki-Modell als FST-Instanziierung]
\label{rem:hu-sawicki-fst}
Das Hu--Sawicki-Modell~\cite{HuSawicki2007}
\[
f(R) = R - \mu^2 \frac{c_1 (R/\mu^2)^n}{1 + c_2 (R/\mu^2)^n}, \quad \mu^2 = H_0^2 \Omega_m, \quad \frac{c_1}{c_2} = \frac{6\Omega_\Lambda}{\Omega_m},
\]
liefert eine konkrete Realisierung des FST-Programms f\"ur Dunkle Energie. Die Modellparameter bilden sich wie folgt auf FST-Gr\"o\ss en ab:
\begin{center}
\begin{tabular}{l|l}
\textbf{Hu--Sawicki} & \textbf{FST-Funktional $\Phi[\rho]$} \\
\hline
$f_{R0} = -n\, c_1 / c_2^2\, (H_0^2 \Omega_m / R_0)^{n+1}$ & Kopplungsst\"arke $\gamma_{\mathrm{eff}}$ \\
Skalaronmasse $m_s^2 = (3 f_{RR})^{-1}$ & Kr\"ummung von $\Phi$ am Minimum \\
Compton-Wellenl\"ange $\lambda_C = 2\pi/m_s$ & \"Ubergangsskala $k_{\mathrm{cross}}$ \\
$n$ (Steilheitsparameter) & Rate des Chameleon-Screenings \\
\end{tabular}
\end{center}
F\"ur $n = 1$ und $|f_{R0}| = 10^{-6}$ betr\"agt die Skalaronmasse bei kosmologischer Dichte $m_s \sim 10^{-33}$~eV ($\lambda_C \sim 30$~Mpc), w\"ahrend sie bei Sonnensystem-Dichte $m_s \sim 10^{-3}$~eV ($\lambda_C \sim 0{,}2$~mm) erreicht, was effizientes Chameleon-Screening liefert. Der gravitative Slip ist $\eta = 1 + 1/(3 f_{RR}\, k^2/a^2 + 1)$, was $\eta \to 4/3$ f\"ur $k \gg a\,m_s$ und $\eta \to 1$ f\"ur $k \ll a\,m_s$ ergibt.

\textbf{Wichtig:} Die FST-Vorhersage $\eta = 5/4$ (aus der allgemeinen $R + \gamma R^{2}$-Analyse in den Gleichungen~\eqref{eq:Poisson-Phi}--\eqref{eq:grav-slip}) und die Hu--Sawicki-Vorhersage $\eta \to 4/3$ unterscheiden sich. Diese Diskrepanz entsteht, weil die quadratische N\"aherung $\gamma R^2$ und die rationale Hu--Sawicki-Funktion verschiedene sub-Horizont-Limiten ergeben: Der $R + \gamma R^2$-Lagrangian liefert Poissongleichungs-Koeffizienten $1/3$ und $2/3$, w\"ahrend das volle nichtlineare Hu--Sawicki-$f(R)$ die Koeffizienten $1/3$ und $1/3$ in den \"aquivalenten Ausdr\"ucken ergibt, was zu $\eta \to (1 + 1/3)/(1 + 0) = 4/3$ f\"uhrt. Die Aufl\"osung erfordert die Berechnung von $\eta$ f\"ur das volle nichtlineare $f(R)$, was modellabh\"angig ist. Die Vorhersage $\eta = 5/4$ ist daher als auf das \emph{quadratische} Regime anwendbar zu verstehen; der pr\"azise Wert f\"ur Hu--Sawicki ist $\eta = 4/3$.
\end{remark}

\subsection{Konkurrierende Ans\"atze und aktuelle Entwicklungen}

Mehrere unabh\"angige Forschungslinien best\"atigen oder hinterfragen die
hier vorgestellte thermodynamische Herleitung der kosmologischen Konstante.

\emph{Holographische Entropie und MOND.}\; Eine aktuelle Analyse in Physics
of the Dark Universe~\cite{MONDEntropy2025} leitet modifizierte
Friedmann-Gleichungen aus verallgemeinerten Entropie-Fl\"achen-Relationen
in Kombination mit dem erweiterten Unsch\"arfeprinzip (EUP) her und
etabliert eine direkte mathematische Verbindung zwischen der
MOND-Beschleunigungsskala und holographischen Entropieschranken. Dies
best\"atigt unabh\"angig, dass holographische thermodynamische Argumente
die makroskopische Gravitationsdynamik in einer Weise modifizieren
k\"onnen, die mit unserer Gleichung~\eqref{eq:a0-MOND} konsistent ist.

\emph{Finsler-Gravitation.}\; Hohmann et al.~\cite{FinslerJCAP2025} zeigen,
dass kinematische Gase, die sich in einer Finsler-Raumzeit (nicht-Riemannsch)
ausbreiten, auf nat\"urliche Weise modifizierte Friedmann-Gleichungen
erzeugen, die beschleunigte Expansion ohne jegliche Dunkle-Energie-Komponente
oder Skalarfeld vorhersagen. Diese rein geometrische Erkl\"arung stellt einen
starken methodischen Konkurrenten zum Skalar-Tensor-Mechanismus des CRM dar:
Sie erreicht dieselbe Ph\"anomenologie allein durch Raumzeit-Geometrie, ohne
thermodynamische oder dissipative Selektionsprinzipien. Ein
Schl\"usseldiskriminant zwischen den beiden Ans\"atzen ist der gravitative
Slip-Parameter $\eta$: Das CRM sagt $\eta = 5/4$ auf sub-Compton-Skalen
voraus, w\"ahrend reine Finsler-Modifikationen generisch
$\eta \neq 5/4$ vorhersagen.

\emph{Informationelle Persistenz.}\; Hunt~\cite{HuntMOND2026} fasst die
Dunkle-Materie-Ph\"anomenologie als \"Uberleben von
Kr\"ummungsinformation unter Vergr\"oberung auf und leitet MOND-artige
Dynamik aus nichtlokaler Gravitation her, ohne Teilchen-Dunkle-Materie zu
ben\"otigen. Dies st\"utzt die Interpretation der ,,informationellen
Persistenz`` von Gleichung~\eqref{eq:a0-MOND} und ordnet das CRM in eine
breitere Familie informationstheoretischer Ans\"atze zur modifizierten
Gravitation ein.

\subsection{Offene Richtungen}

\begin{remark}[Kosmologische Phasenuebergangskette]
Die kosmologische Geschichte laesst sich als Sequenz von Nash-Phasenuebergaengen im Funktional $\Phi[\rho]$ lesen: Strahlungsdominanz $\to$ Materiedominanz $\to$ Dunkle-Energie-Dominanz, wobei jeder einem Wechsel des dominanten Terms der freien Energie entspricht.
\end{remark}

\begin{remark}[Legendre--Fenchel-Dualitaet fuer Dunkle Energie]
Die Legendre--Fenchel-Dualitaet $S \geq F$ liefert eine unabhaengige Konsistenzpruefung: Die thermodynamische Entropie des de-Sitter-Horizonts muss die freie Energie der Vakuumkonfiguration uebersteigen, was $\Lambda_{\mathrm{eff}}$ nach oben beschraenkt.
\end{remark}

\begin{remark}[England-Ungleichung fuer Expansionsrate]
Ein Analogon von Englands dissipativer Adaptations-Ungleichung koennte die kosmische Expansionsrate beschraenken: $\dot{a}/a \geq \Delta S / (k_B \, \Delta t)$, wodurch der Hubble-Parameter mit der Entropieproduktionsrate am kosmologischen Horizont verknuepft wird.
\end{remark}

\begin{remark}[Topologische Ladung als IR-Constraint]
\label{rem:topological-charge}
Das Funktional $\Phi[\rho]$ l\"asst sich zu $\Phi[\rho, Q]$ erweitern, wobei $Q$ eine topologische Ladung (Hopf-Invariante) ist, die die Feldkonfiguration charakterisiert. Minimierung bei festem $Q$ liefert sektorabh\"angige Vakuumenergien $\Lambda_{\mathrm{eff}}(Q)$, wobei das physikalische Vakuum dem Fall $Q = 0$ (triviale Topologie) entspricht. Nicht-triviale Sektoren ($Q \neq 0$) tragen zur Zustandssumme mit einer Entropieunterdr\"uckung $e^{-c|Q|^{3/4}}$ bei (Faddeev--Skyrme-Schranke~\cite{LinYang2004}), was eine nat\"urliche Hierarchie zwischen dem kosmologischen Vakuum und topologischen Defekten liefert. Dies verbindet sich mit der Pattern-A/B-Klassifikation im vereinheitlichten Framework.
\end{remark}

\begin{remark}[Hopfion-Energieschranken als R\"uckwirkung]
\label{rem:hopfion-backreaction}
Die Faddeev--Skyrme-Energieschranke $E \geq c |Q|^{3/4}$ f\"ur verknotete Solitonen~\cite{LinYang2004} liefert eine nicht-quadratische Alternative zur gravitativen Selbstwechselwirkung $V_{\mathrm{grav}} \sim G\rho^2/k^2$. Falls das Dunkle-Energie-Vakuum topologische Anregungen (verknotete Skalarfeldkonfigurationen) tr\"agt, erh\"alt die R\"uckwirkungskostenfunktion einen topologischen Beitrag: $V_{\mathrm{eff}}[\rho] = V_{\mathrm{grav}}[\rho] + \lambda |Q[\rho]|^{3/4}$. Dies ist spekulativ, w\"urde aber einen topologischen Ursprung f\"ur den IR-Cutoff liefern, erg\"anzend zu den dynamischen Abschirmungsmechanismen.
\end{remark}

\begin{remark}[MOND aus verallgemeinerten Entropie-Fl\"achen-Relationen]
\label{rem:mond-entropy}
Jalalzadeh et al.~\cite{MONDEntropy2025} leiten modifizierte Friedmann-Gleichungen aus verallgemeinerten Entropie-Fl\"achen-Relationen (erweitertes Unsch\"arfeprinzip) her. Ihr Ansatz liefert MOND-artige Modifikationen bei niedrigen Beschleunigungen $a < a_0 \sim H_0 c$. Das FST-Funktional $\Phi[\rho]$ teilt den thermodynamischen Ausgangspunkt, unterscheidet sich aber im Mechanismus: $\Phi$ minimiert die freie Energie mit Stabilit\"atsconstraints, w\"ahrend der Entropie-Fl\"achen-Ansatz die Bekenstein--Hawking-Entropie direkt modifiziert. Eine Vereinheitlichung w\"urde erfordern, dass der FST-Stabilit\"atsconstraint (RFEP) \"aquivalent zu einer modifizierten Entropie-Fl\"achen-Relation auf kosmologischen Skalen ist. Dies bleibt eine offene Frage.
\end{remark}

\begin{remark}[Beziehung zur Dynamic Vacuum Field Theory]
\label{rem:dvft}
Die Dynamic Vacuum Field Theory (DVFT) modelliert das Vakuum als dynamisches Skalarfeld und leitet MOND-artige Ph\"anomenologie aus Vakuumfluktuationen her. Die wesentlichen Gemeinsamkeiten und Unterschiede zum FST-Ansatz:
\begin{itemize}
\item \emph{Gemeinsam:} Beide verwenden ein Skalarfeld zur Vermittlung von Dunkle-Energie-Effekten und sagen Abweichungen von $\Lambda$CDM auf galaktischen Skalen voraus.
\item \emph{Unterschiedlich:} DVFT fehlt ein thermodynamisches/Entropie-Argument; das FST-Funktional $\Phi[\rho]$ wird aus dem RFEP (einem Variationsprinzip) abgeleitet, nicht postuliert. DVFT sagt keinen spezifischen Wert von $\eta$ voraus; FST sagt $\eta = 5/4$ voraus.
\end{itemize}
Der Beobachtungsdiskriminant ist klar: Falls $\eta = 5/4$ durch Euclid best\"atigt wird, w\"are DVFT (das generisch $\eta \neq 5/4$ vorhersagt) gegen\"uber FST benachteiligt.
\end{remark}

\begin{remark}[Holographisches Bulk-Boundary-Matching und dichtekorrelierte Constraints]
\label{rem:holographic-matching}
Zwei verwandte Erweiterungen des RG-Rahmenwerks:
\begin{enumerate}
\item \emph{Holographisches Matching:} In einem AdS/CFT-inspirierten Bild ist die UV-Divergenz $\Lambda_0 \sim M_{\mathrm{Pl}}^4$ eine Bulk-Gr\"o\ss e, w\"ahrend $\Lambda_{\mathrm{eff}} \sim H_0^2$ eine Rand-Gr\"o\ss e (Horizont) ist. Das FST-Funktional bildet Bulk-Freiheitsgrade \"uber das RFEP auf den Rand ab, wobei das Skalaron die holographische Projektion vermittelt.
\item \emph{Dichtekorrelierter Constraint:} Die Skalaronmasse $m_s$ kann nichtlinear an die Spur $T^\mu{}_\mu$ des Energie-Impuls-Tensors koppeln: $m_s^2(x) = m_{s,0}^2 (1 + \alpha |T^\mu{}_\mu| / M_{\mathrm{Pl}}^2)$. In dichten Umgebungen ($|T^\mu{}_\mu| \gg M_{\mathrm{Pl}}^2 H_0^2$) divergiert $m_s$ und das Skalaron entkoppelt, was einen weiteren Abschirmungsweg liefert, komplement\"ar zu Cham\"aleon und Vainshtein.
\end{enumerate}
\end{remark}

\begin{remark}[Beziehung zur asymptotischen Sicherheit -- Wetterich 2024]
\label{rem:asymptotic-safety-wetterich}
Wetterich~\cite{Wetterich2024} leitet die Evolution der Dunklen Energie aus einem asymptotisch sicheren UV-Fixpunkt in der Quantengravitation her. Das ,,Cosmon``-Skalarfeld entsteht aus der Skalierungsl\"osung der Renormierungsgruppengleichungen, wobei die Dunkle-Energie-Dichte durch die gr\"o\ss te intrinsische Massenskala bestimmt wird, die der Fluss vom Fixpunkt erzeugt. Der FST-Ansatz teilt die RG-Perspektive, unterscheidet sich aber strukturell: W\"ahrend asymptotische Sicherheit $\Lambda$ aus der UV-Fixpunktstruktur bestimmt, selektiert das FST-Funktional $\Lambda_{\mathrm{eff}}$ \"uber IR-thermodynamisches Gleichgewicht. Beide sagen langsame Evolution von $\Lambda$ voraus, was sie beobachtungsm\"a\ss ig kompatibel, aber theoretisch verschieden macht.
\end{remark}

\begin{remark}[Beziehung zur Covariant Canonical Gauge Gravity]
\label{rem:ccgg}
Das CCGG-Framework von Vasak, Struckmeier und Kirsch~\cite{VasakCCGG2020} leitet Dunkle Energie und Inflation aus lokal kontortierter Raumzeit innerhalb einer De-Donder--Weyl-Hamilton-Formulierung her. Die kosmologische Konstante in CCGG repr\"asentiert die Energie des Raumzeit-Kontinuums, deformiert von seinem (A)dS-Grundzustand zu flacher Geometrie, und entkoppelt sie dadurch von der Vakuumenergie. Dies entsch\"arft das Problem der kosmologischen Konstante auf einem anderen Weg als FST: CCGG verwendet torsionsinduzierte Geometrie, w\"ahrend FST thermodynamische Selektion nutzt. Beide Ans\"atze sagen ein zeitabh\"angiges effektives $\Lambda$ voraus; ihre Beobachtungssignaturen unterscheiden sich im gravitativen Slip ($\eta = 5/4$ f\"ur FST vs.\ torsionsabh\"angige Korrekturen f\"ur CCGG).
\end{remark}

\begin{remark}[Beobachtungsvorhersagen f\"ur $\eta = 5/4$]
\label{rem:forecasts}
Der gravitative Slip-Parameter $\eta = \Phi/\Psi$ liefert das sch\"arfste Unterscheidungsmerkmal zwischen der FST-Vorhersage ($\eta = 5/4$ auf sub-Horizont-Skalen $k \gg a m_s$) und $\Lambda$CDM ($\eta = 1$). Aktuelle und geplante Durchmusterungen bieten folgende Empfindlichkeiten:
\begin{itemize}
\item \emph{Euclid} (DR1 erwartet 2027): $\sigma(\eta) \approx 0.03$--$0.05$ aus schwachem Linseneffekt $+$ Galaxien-Clustering, ausreichend um $\eta = 5/4$ von $\eta = 1$ bei $> 5\sigma$ zu unterscheiden.
\item \emph{DESI} (DR1 2025): $\sigma(\eta) \approx 0.08$ aus BAO $+$ RSD, marginale ($\sim 3\sigma$) Unterscheidung.
\item \emph{LISA} (2030er): Gravitationswellen-Standardsirenen liefern eine unabh\"angige $H_0$-Messung mit $\sigma(H_0)/H_0 \sim 1\%$, was die Hubble-Spannung einschr\"ankt und indirekt $\eta$ testet.
\end{itemize}
Ein Nachweis von $\eta > 1{,}1$ bei $> 3\sigma$ w\"are ein starker Hinweis auf modifizierte Gravitation; $\eta = 1{,}25 \pm 0{,}05$ w\"are ein ,,Smoking Gun`` f\"ur das FST/$f(R)$-Framework.
\end{remark}

\begin{remark}[Fehlerbudget und systematische Unsicherheiten]
\label{rem:error-budget}
Die theoretische Vorhersage $\eta = 5/4$ unterliegt mehreren systematischen Unsicherheiten:
\begin{enumerate}
\item \emph{Standardmodell-Spektrum:} Der Teilcheninhalt bestimmt das Ein-Schleifen-Laufen von $\Lambda$. BSM-Physik (z.\,B.\ zus\"atzliche Skalare) w\"urde $\eta$ um $O(n_{\mathrm{BSM}}/n_{\mathrm{SM}}) \sim 1\%$ verschieben.
\item \emph{Sub-Compton-Regime:} F\"ur $k < a m_s$ bricht die Vorhersage $\eta = 5/4$ zusammen und $\eta \to 1$ (GR-Wiederherstellung). Die \"Ubergangsskala $k_{\mathrm{cross}} = a m_s$ h\"angt von $\gamma$ ab und ist derzeit um einen Faktor $\sim 2$ unsicher.
\item \emph{Renormierungsschema:} Der genaue Wert von $\gamma_0$ ist schema-abh\"angig auf dem $O(10\%)$-Niveau, aber $\eta = 5/4$ ist schema-unabh\"angig (es folgt allein aus der $f(R) = R + \gamma R^2$-Struktur, nicht aus dem Wert von $\gamma$).
\end{enumerate}
\end{remark}

\begin{remark}[Thermodynamische Holographie]
\label{rem:thermo-holography}
Die UV-Divergenz der kosmologischen Konstante ($\Lambda_0 \sim M_{\mathrm{Pl}}^4$) laesst sich holographisch uminterpretieren: Der $M_{\mathrm{Pl}}^4$-Beitrag zaehlt Bulk-Freiheitsgrade, waehrend das physikalische $\Lambda_{\mathrm{eff}} \sim H_0^2$ Rand-Freiheitsgrade auf dem de-Sitter-Horizont zaehlt (Bekenstein--Hawking-Entropie $S = A/(4G) \sim H_0^{-2}/G$). Das FST-Funktional $\Phi[\rho]$ liefert die Bulk-zu-Rand-Abbildung: Die Minimierung von $\Phi$ ueber Bulk-Konfigurationen projiziert auf das rand-dominierte Vakuum und absorbiert die UV-Divergenz in die Renormierung von $\rho_0$.
\end{remark}

\begin{remark}[Asymptotische Sicherheit und der nicht-Gausssche UV-Fixpunkt]
\label{rem:asymptotic-safety-uv}
Im Programm der asymptotischen Sicherheit (Reuter, Wetterich) besitzt der gravitationale RG-Fluss einen nicht-Gaussschen UV-Fixpunkt mit $\Lambda^* / k^2 = \mathrm{const}$ und $G^* k^2 = \mathrm{const}$. Der IR-Fluss von diesem Fixpunkt bestimmt $\Lambda_{\mathrm{eff}}$ bei $k = H_0$. Der FST-Strafterm $V_{\mathrm{grav}} \sim G \rho^2 / k^2$ laesst sich mit der laufenden kosmologischen Konstante bei Skala $k$ identifizieren: $\Lambda(k) = 8\pi G(k) \rho^*(k)$. Falls asymptotische Sicherheit gilt, erbt das FST-Funktional eine natuerliche UV-Vervollstaendigung, und der grosse Wert $\gamma_0 \sim 10^{60}$ erklaert sich durch das RG-Laufen von der Planck-Skala zu $H_0$.
\end{remark}

\begin{remark}[Vainshtein-Abschirmung ueber Galileon-Selbstwechselwirkungen]
\label{rem:vainshtein}
Das Skalaron-Feld $\phi$ im P\"oschl--Teller-Saettigungsmodell kann mit kinetischen Galileon-Selbstwechselwirkungen $\mathcal{L}_{\mathrm{Gal}} = (\partial\phi)^2 \Box\phi / \Lambda_3^3$ ausgestattet werden, wobei $\Lambda_3 = (M_{\mathrm{Pl}} H_0^2)^{1/3}$ die Vainshtein-Skala ist. Innerhalb des Vainshtein-Radius $r_V = (M / (4\pi M_{\mathrm{Pl}} \Lambda_3^3))^{1/3}$ eines massiven Koerpers $M$ dominieren die kinetischen Terme und das Skalaron entkoppelt von der Materie, wodurch GR wiederhergestellt wird. Dies liefert einen komplementaeren Abschirmungsmechanismus zum Chamaeleon-Effekt (der auf dem Anwachsen von $m_s(R)$ beruht) und koennte notwendig sein, falls das P\"oschl--Teller-Potential allein die fuenfte Kraft auf Sonnensystem-Skalen nicht ausreichend unterdrueckt.
\end{remark}

\begin{remark}[Topologische Abschirmung ueber Aharonov--Bohm-Phasen]
\label{rem:topological-screening}
Bei starken baryonischen Dichtegradienten (z.\,B.\ Sterninneren) kann das Skalaron-Feld eine nicht-triviale Windungszahl erwerben, analog zu Aharonov--Bohm-Phasen in der Eichtheorie. Falls $\oint \nabla\phi \cdot d\ell = 2\pi n$ um ein dichtes Objekt gilt, wird die fuenfte Kraft ausserhalb durch die topologische Barriere exponentiell unterdrueckt, was einen von der Skalaron-Masse unabhaengigen Abschirmungsmechanismus liefert. Diese spekulative Moeglichkeit verbindet das FST-Framework mit topologischen Defekten in der Kosmologie (Vilenkin--Shellard) und verdient weitere Untersuchung.
\end{remark}

\begin{remark}[Topologisches RG-Matching und geometrische Phasenuebergaenge]
\label{rem:topological-rg}
Zwei spekulative Erweiterungen des RG-Rahmens:
\begin{enumerate}
\item \emph{Topologisches RG-Matching:} Die quartische UV-Divergenz $\Lambda_0 \sim M_{\mathrm{Pl}}^4$ koennte eine topologische Invariante sein (analog zum Atiyah--Singer-Index), die die Anzahl fermionischer Nullmoden im Gravitationshintergrund zaehlt. Falls ja, ist das UV-IR-Matching kein dynamisches Laufen, sondern eine topologische Identitaet, und $\Lambda_{\mathrm{eff}}$ wird durch die Topologie der de-Sitter-Mannigfaltigkeit bestimmt.
\item \emph{Geometrische Phasenuebergaenge:} Der Uebergang vom kosmologischen Regime ($R \sim H_0^2$) zum Sonnensystem-Regime ($R \sim R_\odot \gg H_0^2$) laesst sich als spontane Symmetriebrechung im FST-Funktional auffassen: Das kosmologische de-Sitter-Vakuum bricht die $\gamma(R)$-Symmetrie, und das Sonnensystem-Vakuum ist eine ,,Higgs-Phase`` mit massivem Skalaron ($m_s \gg H_0$).
\end{enumerate}
\end{remark}

% ============================================================
\subsection*{Danksagung}

\textcolor{gray}{\textit{[Wird in der vollst\"andigen Version erg\"anzt.]}}

% ============================================================
\begin{remark}[Pattern B (Flow-Selektion) Master-Stabilitaet -- problemuebergreifende Struktur]
Diese Arbeit instanziiert Pattern B (Flow-Selektion) aus dem vereinheitlichten Framework \cite{FST-Unified2026}: Strikte Konvexitaet der renormierten freien Energie $F$, kombiniert mit einem neutralen fuehrenden Resolventenzweig und Dominanz zweiter Ordnung, impliziert eine Spektralluecke $\Delta > 0$ fuer den zugehoerigen Transferoperator. Pattern B teilt die variationelle Struktur von Pattern A, ersetzt jedoch den statischen Minimierer durch einen RG-Fluss zu einem Fixpunkt. Die spezifische Instanziierung im vorliegenden Kontext ist: $\Delta = $ Kruemmung von $\Phi[\rho]$ am de-Sitter-Vakuum. Dunkle Energie instanziiert somit Pattern B, das die Positivitaetsarchitektur von Pattern A erbt und dabei das Vakuum ueber einen dynamischen Fluss statt eines statischen Energieminimums selektiert.
\end{remark}

% ===================================================================
% KI-Offenlegung
% ===================================================================
\subsection*{KI-Offenlegung}
Bei der Erstellung dieses Manuskripts wurden KI-gest\"utzte Werkzeuge
(Claude, Anthropic) in unterst\"utzender Funktion eingesetzt,
insbesondere f\"ur Literaturrecherche, Textstrukturierung und
Formulierungshilfe. Der konzeptionelle Entwurf, die wissenschaftliche
Argumentation und die Verantwortung f\"ur den gesamten Inhalt liegen
ausschlie\ss lich beim Autor.

\begin{thebibliography}{99}

\bibitem{Weinberg1989}
S.~Weinberg,
\textit{The cosmological constant problem},
Rev.\ Mod.\ Phys.\ \textbf{61}, 1--23 (1989).

\bibitem{Zeldovich1968}
Ya.~B.~Zeldovich,
\textit{The cosmological constant and the theory of elementary particles},
Sov.\ Phys.\ Usp.\ \textbf{11}, 381--393 (1968).

\bibitem{PeeblesRatra2003}
P.~J.~E.~Peebles and B.~Ratra,
\textit{The cosmological constant and dark energy},
Rev.\ Mod.\ Phys.\ \textbf{75}, 559--606 (2003).

\bibitem{Planck2018}
Planck Collaboration (N.~Aghanim et al.),
\textit{Planck 2018 results. VI. Cosmological parameters},
Astron.\ Astrophys.\ \textbf{641}, A6 (2020).
\href{https://doi.org/10.1051/0004-6361/201833910}{doi:10.1051/0004-6361/201833910}.

\bibitem{Carroll2001}
S.~M.~Carroll,
\textit{The cosmological constant},
Living Rev.\ Relativ.\ \textbf{4}, 1 (2001).
\href{https://doi.org/10.12942/lrr-2001-1}{doi:10.12942/lrr-2001-1}.

\bibitem{Riess1998}
A.~G.~Riess et al.,
\textit{Observational evidence from supernovae for an accelerating
universe and a cosmological constant},
Astron.\ J.\ \textbf{116}, 1009--1038 (1998).

\bibitem{Perlmutter1999}
S.~Perlmutter et al.,
\textit{Measurements of $\Omega$ and $\Lambda$ from 42 high-redshift
supernovae},
Astrophys.\ J.\ \textbf{517}, 565--586 (1999).

\bibitem{MartyushevSeleznev2006}
L.~M.~Martyushev and V.~D.~Seleznev,
\textit{Maximum entropy production principle in physics, chemistry and biology},
Phys.\ Rep.\ \textbf{426}, 1--45 (2006).
\href{https://doi.org/10.1016/j.physrep.2005.12.001}{doi:10.1016/j.physrep.2005.12.001}.

\bibitem{Starobinsky1980}
A.~A.~Starobinsky,
\textit{A new type of isotropic cosmological models without singularity},
Phys.\ Lett.\ B \textbf{91}, 99--102 (1980).
\href{https://doi.org/10.1016/0370-2693(80)90670-X}{doi:10.1016/0370-2693(80)90670-X}.

\bibitem{DESI2025}
DESI Collaboration (A.~G.~Adame et al.),
\textit{DESI DR2 Results II: Measurements of Baryon Acoustic Oscillations and Cosmological Constraints},
arXiv:2503.14738 [astro-ph.CO] (2025).

\bibitem{HuSawicki2007}
W.~Hu and I.~Sawicki,
\textit{Models of $f(R)$ cosmic acceleration that evade solar-system tests},
Phys.\ Rev.\ D \textbf{76}, 064004 (2007).
\href{https://doi.org/10.1103/PhysRevD.76.064004}{doi:10.1103/PhysRevD.76.064004}.

\bibitem{SotiriouFaraoni2010}
T.~P.~Sotiriou and V.~Faraoni,
\textit{$f(R)$ theories of gravity},
Rev.\ Mod.\ Phys.\ \textbf{82}, 451--497 (2010).
\href{https://doi.org/10.1103/RevModPhys.82.451}{doi:10.1103/RevModPhys.82.451}.

\bibitem{BirrellDavies1982}
N.~D.~Birrell and P.~C.~W.~Davies,
\textit{Quantum Fields in Curved Space},
Cambridge Monographs on Mathematical Physics, Cambridge University Press, 1982.
\href{https://doi.org/10.1017/CBO9780511622632}{doi:10.1017/CBO9780511622632}.

\bibitem{MONDEntropy2025}
S.~Jalalzadeh et al.,
\textit{From MOND entropy to extended uncertainty principles},
Phys.\ Dark Universe (2025/2026).

\bibitem{FinslerJCAP2025}
M.~Hohmann et al.,
\textit{Cosmological Finsler spacetimes and accelerated expansion},
JCAP (2025/2026).

\bibitem{HuntMOND2026}
T.~Hunt,
\textit{Dark Matter Phenomenology as Informational Persistence: Nonlocal Gravity, MOND, and the Survival of Curvature Under Coarse-Graining},
ResearchGate preprint, 2026.

\bibitem{GeigerRH2026c}
L.~Geiger, \emph{The Riemann Hypothesis Part III: The Analytical Rank-Finite Certificate}, Zenodo preprint, March 2026.
\href{https://doi.org/10.5281/zenodo.19035845}{doi:10.5281/zenodo.19035845}.

\bibitem{FST-Unified2026}
L.~Geiger, \emph{The Renormalized Free Energy Framework: A Unified Resolution of Structural Bottlenecks}, Zenodo preprint, 2026.
\href{https://doi.org/10.5281/zenodo.19036191}{doi:10.5281/zenodo.19036191}.

\bibitem{FST-NS2026}
L.~Geiger, \emph{A Thermodynamic Obstruction to Finite-Time Blow-Up in the 3D Navier--Stokes Equations (Conditional Framework)},
Preprint, 2026.

\bibitem{Geiger2026-YM}
L.~Geiger, \emph{Volume-Independent Spectral Gap for Strong-Coupling Lattice Gauge Theory via Poincar\'e--Dobrushin Machinery},
Preprint, 2026.

\bibitem{Geiger2026-BSD}
L.~Geiger, \emph{The BSD Conjecture as a Positivity Normal-Form Theorem: Height Saturation and the Impossibility of Rank-Order Discrepancy},
Preprint, 2026.

\bibitem{CopelandLiddleWands1998}
E.~J.~Copeland, A.~R.~Liddle, and D.~Wands,
\textit{Exponential quintessence},
Phys.\ Rev.\ D \textbf{57}, 4686--4690 (1998).
\href{https://doi.org/10.1103/PhysRevD.57.4686}{doi:10.1103/PhysRevD.57.4686}.

\bibitem{Li2004}
M.~Li,
\textit{A model of holographic dark energy},
Phys.\ Lett.\ B \textbf{603}, 1--5 (2004).
\href{https://doi.org/10.1016/j.physletb.2004.10.014}{doi:10.1016/j.physletb.2004.10.014}.

\bibitem{BoussoPolchinski2000}
R.~Bousso and J.~Polchinski,
\textit{Quantization of four-form fluxes and dynamical neutralization of the cosmological constant},
JHEP \textbf{0006}, 006 (2000).
\href{https://doi.org/10.1088/1126-6708/2000/06/006}{doi:10.1088/1126-6708/2000/06/006}.

\bibitem{Susskind2003}
L.~Susskind,
\textit{The anthropic landscape of string theory},
in: \textit{Universe or Multiverse?}, Cambridge University Press, 2003.
\href{https://doi.org/10.1017/CBO9781107050990.018}{doi:10.1017/CBO9781107050990.018}.

\bibitem{KaloperPadilla2014}
N.~Kaloper and A.~Padilla,
\textit{Sequestering the standard model vacuum energy},
Phys.\ Rev.\ Lett.\ \textbf{112}, 091304 (2014).
\href{https://doi.org/10.1103/PhysRevLett.112.091304}{doi:10.1103/PhysRevLett.112.091304}.

\bibitem{LinYang2004}
F.~Lin and Y.~Yang,
\textit{Existence of energy minimizers as stable knotted solitons in the Faddeev model},
Commun.\ Math.\ Phys.\ \textbf{249}, 273--303 (2004).
\href{https://doi.org/10.1007/s00220-004-1110-y}{doi:10.1007/s00220-004-1110-y}.

\bibitem{Wetterich2024}
C.~Wetterich,
\textit{Dark energy evolution from quantum gravity},
arXiv:2407.03465 [hep-th] (2024).

\bibitem{VasakCCGG2020}
D.~Vasak, J.~Struckmeier, and J.~Kirsch,
\textit{Dark energy and inflation invoked in CCGG by locally contorted space-time},
Eur.\ Phys.\ J.\ Plus \textbf{135}, 404 (2020).
\href{https://doi.org/10.1140/epjp/s13360-020-00415-7}{doi:10.1140/epjp/s13360-020-00415-7}.

\end{thebibliography}

\end{document}
